\documentclass[11pt]{article}
\usepackage[a4paper,margin=28mm]{geometry}
\usepackage{amsmath,amssymb,amsthm,booktabs,array}
\usepackage[colorlinks=true,linkcolor=blue,citecolor=blue,urlcolor=blue]{hyperref}
\newtheorem{theorem}{Theorem}[section]
\newtheorem{proposition}[theorem]{Proposition}
\newtheorem{lemma}[theorem]{Lemma}
\newtheorem{corollary}[theorem]{Corollary}
\theoremstyle{remark}
\newtheorem{remark}[theorem]{Remark}
\numberwithin{equation}{section}
\setlength{\emergencystretch}{2em}
\title{Small-support inequalities for reflected max-convolution}
\author{Benjamin Guo\\Hraness (\href{https://hraness.com}{\texttt{hraness.com}})\thanks{Developed with AI research agents. The hand proofs use the cited classical small-sumset theorem; the arbitrary-weight five-point result also uses rational certificates checked by the supplied program.}}
\date{30 September 2026}
\hypersetup{pdftitle={Small-support inequalities for reflected max-convolution},pdfauthor={Benjamin Guo}}

\begin{document}
\maketitle
\begin{abstract}
For a nonnegative kernel on at most four integers, we bound its weighted
squared mass under a coefficientwise ordinary-autoconvolution constraint by
the fourth root of its fourfold max-convolution mass. Finite linear
programming duality identifies the constrained maximum with a universal
translated-kernel covering constant. The bound therefore gives a reflected
fourfold max-convolution inequality, strict for positive kernels on four
points and nonzero inputs. We classify the five pair-sum collision patterns on four points and
prove the corresponding sharp unweighted covering constants. The
five-point Sidon set $\{0,2,7,8,11\}$ violates the stronger covering bound by
$1/16$, establishing the support-size limit of that bound. Nevertheless,
the reflected inequality holds for every five-point integer support when
the kernel is constant, with arbitrary nonnegative input and denominator
weights. The proof combines an energy estimate, a bound using the range
of the denominator weights, and a small-sumset argument. For the specific
support $\{0,2,7,8,11\}$, we also prove the reflected inequality for
arbitrary nonnegative kernel weights. A sharper allocation estimate and
240 rational coefficient certificates cover all weight orders, with
strictness when all five kernel weights are positive and $f,g$ are nonzero.
\end{abstract}

\tableofcontents

\section{The inequality and its covering constant}
All functions below are nonnegative and finitely supported on $\mathbb Z$.
We distinguish ordinary convolution from max-convolution:
\[
 (u*v)(s)=\sum_{p+q=s}u(p)v(q),\qquad
 (u\star v)(s)=\max_{p+q=s}u(p)v(q).
\]
The maximum is zero when no pair contributes. We write
$\widetilde u(p)=u(-p)$ and $\|u\|_1=\sum_pu(p)$.
For a finite nonempty set $F$ and a kernel $k\ge0$ supported on $F$, define
\begin{equation}\label{eq:definitions}
 N_F(k)=\max_{\substack{x\ge0,\ \operatorname{supp}x\subseteq F\\x*x\le1}}
             \sum_{p\in F}k(p)x(p)^2,
 \qquad T_F(k)=\|k\star k\star k\star k\|_1.
\end{equation}
The constraint is coefficientwise. Each coordinate of $x$ is at most one,
so the maximum exists by compactness.

\begin{theorem}\label{thm:main}
If $1\le |F|\le4$, then every nonnegative kernel supported on $F$ satisfies
\begin{equation}\label{eq:main}
 N_F(k)^4\le T_F(k).
\end{equation}
If $|F|=4$ and $k$ is positive at every point of $F$, the inequality is strict.
For $F=\{0,2,7,8,11\}$ and $k=1_F$, instead
\[
 N_F(k)^4\ge(5/2)^4=39+1/16>39=T_F(k).
\]
Thus four is the largest support size for which \eqref{eq:main} holds
universally.
\end{theorem}

The theorem concerns a covering bound for the reflected max-convolution
ratio. The connection follows from a general finite-dimensional lemma.

\begin{lemma}[Universal translated covering]\label{lem:cover}
For $g\ge0$, nonzero and supported on $F$, define
\begin{equation}\label{eq:cover}
 L_F(g,k)=\min_{\lambda_t\ge0}
 \left\{\sum_{t\in F+F}\lambda_t:
 \sum_{q\in F}g(q)\lambda_{p+q}\ge g(p)k(p)\quad(p\in F)\right\}.
\end{equation}
Then
\begin{equation}\label{eq:universal}
 \sup_{g\ne0,\ g\ge0}L_F(g,k)=N_F(k).
\end{equation}
The same supremum results if $g$ is required to be positive throughout $F$.
For every nonzero finitely supported $f\ge0$,
\begin{equation}\label{eq:ratio}
 \|f\star\widetilde{gk}\|_1
 \le L_F(g,k)\|f\star g\|_1
 \le N_F(k)\|f\star g\|_1.
\end{equation}
Here $gk$ denotes pointwise multiplication.
\end{lemma}
\begin{proof}
Fix $q_0\in F$ with $g(q_0)>0$. Set
$\lambda_{p+q_0}=g(p)k(p)/g(q_0)$ for $p\in F$, and set all remaining
variables to zero. The shifts $p+q_0$ are distinct, so this is feasible.
A bounded minimizing level set is compact. The finite linear-programming
dual maximizes $\sum_pk(p)g(p)z(p)$ subject to
$z\ge0$ and $z*g\le1$. It contains zero and is bounded, since
$z(p)\le1/g(q_0)$. Both extrema are attained and are equal by finite
linear-programming duality.

For a feasible dual vector set $x(p)=\sqrt{g(p)z(p)}$. Coefficientwise,
\[
 (x*x)(s)\le\frac12\sum_{p+q=s}
       \bigl(g(p)z(q)+g(q)z(p)\bigr)=(z*g)(s)\le1.
\]
Thus $L_F(g,k)\le N_F(k)$. Conversely, for a nonzero maximizer $x$ in
\eqref{eq:definitions}, taking $g=z=x$ attains its objective in the dual.
If $k\ne0$, homogeneity ensures $\|x*x\|_\infty=1$ at a maximizer.
With $g_\epsilon=x+\epsilon1_F$ and
$z_\epsilon=g_\epsilon/\|g_\epsilon*g_\epsilon\|_\infty$, the dual value
tends to $N_F(k)$. This proves the positive-$g$ assertion; $k=0$ is immediate.

An attained cover gives, with $h=\widetilde{gk}$,
\[
 h(v)\le\sum_t\lambda_tg(v+t),\qquad
 (f\star h)(u)\le\sum_t\lambda_t(f\star g)(u+t).
\]
The first inequality is the constraint at $p=-v$ when $v\in-F$ and is
automatic elsewhere. Summing the second over $u$ proves \eqref{eq:ratio}.
\end{proof}

\begin{corollary}\label{cor:h4}
If $|F|\le4$, $g\ge0$ is nonzero and supported on $F$, and $k\ge0$ is
supported on $F$, then every nonzero finitely supported $f\ge0$ satisfies
\[
 \|f\star\widetilde{gk}\|_1^4
 \le \|f\star g\|_1^4\,\|k\star k\star k\star k\|_1.
\]
The inequality is strict when $|F|=4$ and $k>0$ on $F$, including when
some values of $g$ are zero.
\end{corollary}
\begin{proof}
Apply Theorem~\ref{thm:main} to Lemma~\ref{lem:cover};
$\|f\star g\|_1>0$.
\end{proof}

Theorem~\ref{thm:five-point-weighted} proves the same reflected inequality
for arbitrary kernels on $F=\{0,2,7,8,11\}$, where the stronger bound
\eqref{eq:main} fails. Its proof uses the actual common input $f,g$
to improve the covering estimate before comparing with $T_F(k)$.

Translation and nonzero dilation of $F$ relabel pair sums and fourfold sums
injectively and preserve $N_F$ and $T_F$. The translated covering argument
uses shifts in the actual integer support, so it imposes no residue-class
condition on $f$. Reflection is included. Multiplying $k$ by a constant
scales $N_F$ linearly and $T_F$ to the fourth power.

\subsection{Relation to sumset inequalities}
Gyarmati, Hennecart and Ruzsa~\cite[p.~177, equation~(9)]{ghr} asked when
\[
 |A-B|\le |A+B|\,|hB|^{1/h}
\]
holds for all finite nonempty sets in an abelian group, where $hB$ is the
repeated sumset.
Their paper gives a sixfold counterexample and leaves four and five
unresolved~\cite[p.~182]{ghr}. With indicator functions and $k=1_F$,
Corollary~\ref{cor:h4} and Theorem~\ref{thm:five-point} give the fourfold
inequality for $B=F$ of size at most five. Both allow arbitrary $f$ and
$g$; the four-point result also permits arbitrary kernel weights $k$.
The five-point example in Theorem~\ref{thm:main} concerns $N_F$ and
satisfies the reflected inequality for every kernel by
Theorem~\ref{thm:five-point-weighted}.

Max-convolution and weighted sumset parameters are established tools.
Matolcsi, Ruzsa, Shakan and Zhelezov~\cite[Definition~8.1]{mrsz} use the
operation above and relate a weighted induced-tripling parameter to a set
parameter in their Theorem~8.6. The cube norm inequality of Becker,
Ivanisvili, Krachun and Madrid~\cite[Theorem~1.3]{bikm} and the
mixed-alphabet bounds of Hosle and Ivanisvili~\cite[Theorem~1.5]{hi}
concern lower bounds for max-convolution and sumset expansion.
Green, Matolcsi, Ruzsa, Shakan and Zhelezov~\cite[Proposition~2.1]{gmrsz}
prove the weighted discrete Pr\'ekopa--Leindler inequality
\[
 \sum_n\max\{p(a\star b)(n),(1-p)(a\star b)(n-1)\}
 \ge \|a\|_2\|b\|_2\qquad(0\le p\le1),
\]
and deduce a triple-sum lower bound for subsets of quasicubes in their
Theorem~1.1. Their proposition is equivalent to the $p=q=2$ case of
Theorem~11.1 of Matolcsi et al.~\cite{mrsz}.
For ordinary autoconvolution, Cloninger and
Steinerberger~\cite[Section~1.2]{cs} and Ivanisvili and
Xie~\cite[Theorem~5]{ix} study continuous $L^1$ mass on a fixed interval.
That normalization differs from the weighted squared input mass in $N_F(k)$
and its comparison with $T_F(k)$.

These differences in theorem statements establish neither novelty nor
non-implication. The covering proofs are self-contained; the five-point
result also uses Freiman's classical $3k-4$ theorem, in the form recalled
by Bollob\'as, Leader and Tiba~\cite[p.~1]{blt}. Priority in the
wider literature remains unestablished. In particular, the mathematical
comparison with the weighted Pr\'ekopa--Leindler results still requires
control of the linked ratio involving $f\star g$ and
$f\star\widetilde{gk}$, rather than a lower bound for an independent
max-convolution. The companion
\href{https://github.com/hraness/algal-lab/blob/main/papers/sumset-literature.md}
{literature comparison} records the exact source versions, reading limits
and unreproduced numerical dependencies of the continuous results.

\subsection{A coefficient comparison used below}
\label{sec:coefficient}
Let $v\ge0$ be a profile on $F$. Group the expansion of $(k\cdot v)^4$
by the sum of its four support positions. If the total coefficient weight
in each group is at most one, then
\begin{equation}\label{eq:coefficient}
 (k\cdot v)^4\le T_F(k).
\end{equation}
Indeed, each kernel monomial in that group is at most its max-convolution
value. The same argument applies after any stated nonnegative
AM--GM replacement of monomials. Such a replacement may move weight to
another sum; the proof must account for the coefficient total at every
destination. A strict deficit at a group with positive maximum proves
strictness. A profile used in \eqref{eq:coefficient} need not itself be
feasible for the ordinary-autoconvolution constraint.

\section{The five pair-sum patterns}
\label{sec:classification}

The ordinary-convolution constraints depend only on equalities between
unordered pair sums, with repeated summands allowed.  For four integer
points these equalities give five patterns.

\begin{proposition}\label{c-classification}
Every four-point integer set belongs to exactly one of the following
classes: a Sidon set, an isolated three-term arithmetic progression, a
proper parallelogram, two three-term arithmetic progressions, or a
four-term arithmetic progression.  Here a Sidon set has distinct unordered
pair sums.  An isolated progression has exactly one nontrivial pair-sum
equality, of the form $p+r=2q$.  A proper parallelogram has exactly one
nontrivial pair-sum equality, with four distinct summands.

Up to translation, reflection, and integer dilation, the class with two
progressions is represented by $\{0,1,2,4\}$.
\end{proposition}

\begin{proof}
Translate the set to $\{0,a,b,c\}$, where $0<a<b<c$.  A shared summand
can be cancelled from any pair-sum equality.  The remaining nontrivial
possibilities are precisely
\begin{equation}\label{c-relations}
 b=2a,\qquad c=2a,\qquad c=2b,\qquad
 c=2b-a,\qquad c=a+b.
\end{equation}
The first four assert that one point is the midpoint of two others.
The last is the only possible equality between pairs using all four points.

The compatible intersections are as follows.  The equations $b=2a$ and
$c=2b$ give $\{0,a,2a,4a\}$.  The equations $c=2a$ and $c=2b-a$ give
$b=3a/2$, hence a reflected copy of $\{0,d,2d,4d\}$ with $d=a/2$.
Any two of $b=2a$, $c=2b-a$, and $c=a+b$ imply the third and give
$\{0,a,2a,3a\}$.  Every other pair of equations in
\eqref{c-relations} contradicts $0<a<b<c$.  The possibilities are
therefore no equality, one midpoint equality, one cross-pair equality,
the two-progressions case, and the four-term progression.
\end{proof}

This classification concerns pair sums.  A class can contain further
equalities between fourfold sums, which must also be accounted for when
bounding $T_F(k)$.

\section{Sidon supports}
\label{sec:sidon}

\begin{lemma}\label{c-level-set}
For every finitely supported nonnegative function $h$ on $\mathbb Z$,
\begin{equation}\label{c-level-bound}
 \|h\star h\|_1\ge 2\sum_s h(s)^2-\|h\|_\infty^2.
\end{equation}
\end{lemma}

\begin{proof}
The zero function is immediate.  For $0\le t<\|h\|_\infty^2$, put
$A_t=\{s:h(s)^2>t\}$.  Then $A_t$ is nonempty and
$A_t+A_t\subseteq\{s:(h\star h)(s)>t\}$.  If a finite integer set has
ordered elements $a_1<\cdots<a_m$, the chain
\[
 2a_1<a_1+a_2<\cdots<a_1+a_m<a_2+a_m<\cdots<2a_m
\]
shows that $|A+A|\ge2|A|-1$.  Integrating the resulting inequality for
superlevel-set cardinalities over $0\le t<\|h\|_\infty^2$ proves
\eqref{c-level-bound}.
\end{proof}

\begin{proposition}\label{c-sidon}
Let $F$ be a finite nonempty Sidon set and let $k\ge0$ on $F$.  Put
$K=\sum_{p\in F}k(p)$ and $a=\max_{p\in F}k(p)$.  Then
\begin{equation}\label{c-sidon-norm}
 N_F(k)=\max\left\{\frac K2,\frac{K+3a}{4}\right\}.
\end{equation}
If $|F|\le4$, then $N_F(k)^4\le T_F(k)$, with strict inequality
whenever at least two values of $k$ are positive.
\end{proposition}

\begin{proof}
The Sidon property makes the ordinary-convolution constraints equivalent to
\[
 w_p\le1,\qquad w_pw_q\le\frac14\quad(p\ne q),
 \qquad w_p=x(p)^2\ge0.
\]
If all $w_p\le1/2$, the objective is at most $K/2$.  Otherwise let
$w_j=u\in[1/2,1]$ be largest.  Every other coordinate is at most
$1/(4u)$, so
\[
 \sum_p k(p)w_p\le k(j)u+\frac{K-k(j)}{4u}.
\]
This is convex in $u$, and its endpoint values are $K/2$ and
$k(j)+(K-k(j))/4$.  Maximizing over $j$ proves the upper bound in
\eqref{c-sidon-norm}.  The profiles with all coordinates $1/2$, and
with one coordinate $1$ and all others $1/4$, are feasible.  They attain
the stated values.

Suppose now that $|F|\le4$.  The case $k=0$ is immediate.  Apply
Lemma~\ref{c-level-set} to $h=k\star k$.  Distinct unordered pair sums
give
\[
 \sum_s h(s)^2
 =\frac12\left[\left(\sum_p k(p)^2\right)^2+\sum_p k(p)^4\right],
 \qquad \|h\|_\infty=a^2.
\]
Associativity of max-convolution and Cauchy--Schwarz therefore imply
\begin{equation}\label{c-sidon-half}
 T_F(k)\ge
 \left(\sum_p k(p)^2\right)^2+\sum_p k(p)^4-a^4
 \ge\left(\frac K2\right)^4.
\end{equation}
The last comparison is strict if at least two kernel values are positive,
because then $\sum_p k(p)^4-a^4>0$.

For the other term in \eqref{c-sidon-norm}, choose a point $p_0$ with
$k(p_0)=a$ and write $R=K-a$.  Fixing two copies of $p_0$ in the
fourfold maximum gives distinct sums indexed by unordered pairs from $F$.
Thus
\begin{equation}\label{c-pivot-bound}
 T_F(k)\ge a^4+a^3R+a^2h_2
 \ge a^4+a^3R+\frac23a^2R^2,
\end{equation}
where $h_2$ is the sum of the squares and pair products of the other
kernel values.  There are at most three such values, so
$h_2=(R^2+\sum_{p\ne p_0}k(p)^2)/2\ge2R^2/3$.
Since $R\le3a$, expansion gives
\begin{equation}\label{c-quarter-bound}
 \left(a+\frac R4\right)^4
 \le a^4+a^3R+\frac{153}{256}a^2R^2.
\end{equation}
Indeed the coefficient after bounding the cubic and quartic terms is
$3/8+3/16+9/256=153/256$.  The difference between the right sides of
\eqref{c-pivot-bound} and \eqref{c-quarter-bound} is at least
$53a^2R^2/768$, positive when $R>0$.  This proves both comparisons in
\eqref{c-sidon-norm}, including the asserted strictness.
\end{proof}

\section{Proper parallelograms}
\label{sec:parallelogram}

Write
\begin{equation}\label{c-parallelogram-set}
 F=\{0,a,b,a+b\},\qquad 0<a<b,\qquad b\ne2a,
\end{equation}
and index its coordinates in this order.  The opposite pairs are
$\{0,3\}$ and $\{1,2\}$.  Define a collection $\mathcal P$ of sixteen
squared profiles:
\begin{itemize}
 \item Four profiles have three coordinates $1/2$ and one coordinate zero.
 \item Eight profiles have one coordinate $1$, its opposite coordinate
       and one adjacent coordinate $1/4$, and the remaining coordinate zero.
 \item Four profiles have one coordinate $1$, its two adjacent coordinates
       $1/4$, and its opposite coordinate $1/16$.
\end{itemize}
We call the last four profiles tensor profiles, since their entries are
products of the two-coordinate vector $(1,1/4)$.

\begin{proposition}\label{c-parallelogram}
For every $k\ge0$ on \eqref{c-parallelogram-set},
\begin{equation}\label{c-parallelogram-norm}
 N_F(k)=\max_{v\in\mathcal P}k\cdot v,
 \qquad (k\cdot v)^4\le T_F(k)\quad(v\in\mathcal P).
\end{equation}
If $k$ is positive at all four points, every displayed quartic comparison
is strict, and hence $N_F(k)^4<T_F(k)$.
\end{proposition}

\subsection{Reduction to sixteen profiles}

Ordinary-convolution feasibility is equivalent to
\begin{equation}\label{c-parallelogram-feasible}
 x_i^2\le1,\qquad 2x_ix_j\le1\ (i\ne j),\qquad
 x_0x_3+x_1x_2\le\frac12.
\end{equation}
The last constraint is the coefficient at $a+b$; its separate opposite-pair
bounds are redundant.  Each profile in $\mathcal P$ is feasible.
For the first two types the sum of the two opposite products of square
roots is $1/2$; for a tensor profile both products are $1/4$.

Put $w_i=x_i^2$.  First suppose every $w_i\le1/2$.  With
$r=x_0x_3$ and $s=x_1x_2$, the inequality
$u+v\le1/2+2uv$ for $u,v\in[0,1/2]$ gives
\[
 \sum_iw_i\le1+2r^2+2s^2\le1+2(r+s)^2\le\frac32.
\]
Increase coordinates within $[0,1/2]$ until their sum is $3/2$.
The resulting vector lies in the convex hull of the four profiles of the
first type: for such a vector $y$, the coefficient of the profile with its
$i$th coordinate zero is $1-2y_i$.  These coefficients are nonnegative
and sum to one.

Otherwise relabel corners, preserving opposite pairs, so that $w_0=u\ge1/2$,
and put $L=1/(4u)$.  Then $w_1,w_2,w_3\le L$, and
\[
 w_3\le L\left(1-2\sqrt{w_1w_2}\right)^2.
\]
For fixed $u$, the upper bound
\[
 k_0u+k_1w_1+k_2w_2
       +k_3L\left(1-2\sqrt{w_1w_2}\right)^2
\]
is separately convex in $w_1,w_2\in[0,L]$.  Its only nonlinear term in
either variable is a nonpositive multiple of that variable's square root.
It is therefore enough to consider the four corners of this square.

The corner $(0,0)$ is dominated by either $(L,0)$ or $(0,L)$.  These
give $(u,L,0,L)$ and $(u,0,L,L)$, whose objectives are convex in $u$.
Their endpoint values correspond to profiles of the first two types.
The last corner gives
\[
 b(u)=\left(u,L,L,\frac{(2u-1)^2}{16u^3}\right).
\]
For this curve, set
\[
 t=2u-1,\qquad \nu=1-t,\qquad
 s=\frac{(2u-1)(1-u)}{4u},
\]
and define
\[
 U=(1,1/4,1/4,1/16),\quad W_0=(1/2,1/2,1/2,0),
\]
\[
 W_1=(1/2,1/2,0,1/2),\qquad W_2=(1/2,0,1/2,1/2).
\]
The coefficients of
\[
 y=tU+(\nu-4s)W_0+2sW_1+2sW_2
\]
are nonnegative and sum to one, since $4s=t(1-u)/u\le2(1-u)=\nu$.
The first three coordinates of $y$ are $(u,L,L)$ and its last coordinate
is $t/16+2s$.  Moreover,
\[
 \frac{(2u-1)^2}{16u^3}-\frac t{16}
 =\frac{t(1-u)(u^2+u-1)}{16u^3}
 \le\frac s4\le2s.
\]
Thus $y$ dominates $b(u)$.  These reductions prove the upper bound in
\eqref{c-parallelogram-norm}; feasibility of the listed profiles proves
equality.

\subsection{Profiles supported on three corners}

Every three-corner subset of a proper parallelogram is Sidon.  An additional
midpoint equality would force $b=2a$ or $a=b$.  For a kernel restricted to
such a triangle, Proposition~\ref{c-sidon} bounds both its half-sum and every
unit-and-two-quarter form.  These are exactly the objectives of the twelve
profiles with a zero coordinate.  Monotonicity under enlarging the support
then proves their quartic comparisons with $T_F(k)$.  For positive $k$,
their three-point restrictions have at least two positive values, so the
comparisons are strict.

\subsection{Tensor profiles and their collision grid}

Reflection around $(a+b)/2$ reduces the four tensor profiles to a pivot at
$0$ and a pivot at $a$.  Relative to the pivot, write the kernel values
as $A,B,C,D$, with the tensor objective
\[
 E=A+B/4+C/4+D/16.
\]
For the outer pivot the offsets are $0,a,b,a+b$.  For the inner pivot
they are $0,-a,b,b-a$, so that $A=k(a)$, $B=k(0)$, $C=k(a+b)$, and
$D=k(b)$ in that case.

A monomial with multiplicities $n_A,n_B,n_C,n_D$ has grid coordinates
$i=n_B+n_D$ and $j=n_C+n_D$.  Both belong to $\{0,1,2,3,4\}$.
The total coefficient of its grid point in $E^4$ is
\begin{equation}\label{c-grid-weights}
 c_ic_j,\qquad (c_0,c_1,c_2,c_3,c_4)=(1,1,3/8,1/16,1/256),
\end{equation}
because the coefficient-generating polynomial is
\[
 (1+z/4+w/4+zw/16)^4=(1+z/4)^4(1+w/4)^4.
\]
The actual sum is $ia+jb$ for the outer pivot and $-ia+jb$ for the inner
pivot.  Grouping by that sum proves $E^4\le T_F(k)$ whenever every
group's nonnegative coefficient weights add to at most one.  AM--GM
inequalities can first move terms between groups, provided each replacement
is a legal fourfold monomial at its new sum.

Write $b/a=r/s$ in lowest terms, with $r>s>0$.  Two grid points in
$\{0,\ldots,4\}^2$ collide only if $r\le4$.  Excluding $r/s=2$ leaves
\begin{equation}\label{c-exceptional-ratios}
 r/s\in\{3,4,3/2,4/3\}.
\end{equation}
The collision step is $(r,-s)$ for the outer pivot and $(r,s)$ for the
inner pivot.  Since $r\ge3$, no group contains three grid points.
Outside \eqref{c-exceptional-ratios}, \eqref{c-grid-weights} proves the
quartic comparison directly.

We use two elementary inequalities:
\begin{align}
 U+\epsilon V&\le\max(U,V)+\epsilon\min(U,V)
     &&(U,V\ge0,\ 0\le\epsilon\le1),\label{c-max-min}\\
 \min(A^3C,X^4)&\le A^2X^2+A^2C^2
     &&(A,C,X\ge0).\label{c-absorption}
\end{align}
The first follows by ordering $U,V$.  For the second, the case $A=0$
is immediate.  Otherwise divide by $A^4$ and put $x=X/A$, $y=C/A$.
If $x\le1$, use $\min(y,x^4)\le x^4\le x^2$; if $x\ge1$ and
$y\le1$, use $y\le x^2$; if $y\ge1$, use $y\le y^2$.

\subsection{Outer-pivot collisions}

At ratios $3/2$ and $4/3$, every colliding pair has total weight at most
$3/8+1/16=7/16$.  Indeed the point with smaller first coordinate has
second coordinate at least two, and the other point has first coordinate
at least three.  Thus the coefficient comparison applies directly.

At ratio three, the only groups of weight greater than one are
$(0,1)\sim(3,0)$ and $(1,1)\sim(4,0)$.  Their polynomials are
\[
 A^3C+AB^3/16,\qquad
 3A^2BC/4+A^3D/4+B^4/256.
\]
Use $AB^3/16\le(A^2B^2+B^4)/32$ in the first group.  The resulting
total coefficient of $B^4$ in the second group is $\delta=9/256$.
Keep
\[
 (3/4-\delta)A^2BC+A^3D/4+\delta B^4
\]
there; its coefficient sum is one.  The displaced term satisfies
$\delta A^2BC\le\delta(A^2B^2+A^2C^2)/2$.
The extra coefficients assigned to $A^2B^2$ and $A^2C^2$ are
$25/512$ and $9/512$.  Their entire original group weights are
$3/8$ and $7/16$, respectively, so the final weights are
$217/512$ and $233/512$.  Both are below one.

At ratio four, the sole overloaded group is $A^3C+B^4/256$.
Equations~\eqref{c-max-min}--\eqref{c-absorption}, with
$\epsilon=1/256$ and $X=B$, bound it by its maximum plus
$(A^2B^2+A^2C^2)/256$.  The two recipient groups originally have weights
$3/8$ and $97/256$; their final weights are $97/256$ and $98/256$.
Every other group already has weight at most one.

\subsection{Inner-pivot collisions}

An inner-pivot group can be overloaded only if the smaller grid point has
both coordinates at most one.  Otherwise its weight is at most $3/8$,
and that of its partner is at most $1/16$.
For ratios three, four, and $3/2$, remove the high-point polynomials in
the following table and keep their low-point partners in full.
\begin{center}
\begin{tabular}{cl}
\toprule
High grid point & Polynomial \\
\midrule
$(3,1)$ & $B^3C/64+3AB^2D/64$ \\
$(4,1)$ & $B^3D/256$ \\
$(3,2)$ & $3B^2CD/256+3ABD^2/256$ \\
$(4,2)$ & $3B^2D^2/2048$ \\
$(3,3)$ & $3BCD^2/1024+AD^3/1024$ \\
$(4,3)$ & $BD^3/4096$ \\
\bottomrule
\end{tabular}
\end{center}
The entries follow from the multinomial expansion of $E^4$.  The sets
of removed points and their total coefficient weights are
\begin{center}
\begin{tabular}{clc}
\toprule
Ratio & Removed points & Total weight \\
\midrule
$3$ & $(3,1),(4,1),(3,2),(4,2)$ & $187/2048$ \\
$4$ & $(4,1),(4,2)$ & $11/2048$ \\
$3/2$ & $(3,2),(4,2),(3,3),(4,3)$ & $119/4096$ \\
\bottomrule
\end{tabular}
\end{center}
Each total is below $1/8$.  Every unweighted monomial in the first table
is at most
\[
 S=B^4+C^4+D^4+A^2B^2+A^2D^2.
\]
For monomials without $A$, this follows from weighted AM--GM.  The three
remaining forms satisfy
\[
 AB^2D\le\frac{A^2D^2+B^4}{2},\qquad
 ABD^2\le\frac{A^2B^2+D^4}{2},\qquad
 AD^3\le\frac{A^2D^2+D^4}{2}.
\]
The removed polynomial is consequently bounded by $S/8$.

The five recipient monomials occupy distinct groups, outside the removed
groups.  In the order $B^4,C^4,D^4,A^2B^2,A^2D^2$, their original
weights are
\[
 \frac1{256},\quad\frac1{256},\quad
 \begin{cases}
  1/16+1/65536,&r/s=3\text{ or }4,\\
  3/8+1/65536,&r/s=3/2,
 \end{cases}
 \quad\frac38,\quad\frac9{64}.
\]
These are full group weights, including any collision partner.  For example,
the partner of the $D^4$ point $(4,4)$ is $(1,3)$, $(0,3)$, or $(1,2)$
at the respective ratios.  The other recipient grid points are
$(4,0),(0,4),(2,0),(2,2)$ and are unpaired.  Adding $1/8$ to each
weight leaves every total below one, proving all three cases.

At ratio $4/3$, the two overloaded groups are
\[
 A^4+BD^3/4096,\qquad A^3C+D^4/65536.
\]
Weighted AM--GM gives
\[
 BD^3/4096+D^4/65536
 \le B^4/16384+13D^4/65536.
\]
Keep $A^4$ in its group and combine the $D^4$ term with $A^3C$ using
\eqref{c-max-min}--\eqref{c-absorption}, with
$\epsilon=13/65536$ and $X=D$.  This adds $1/16384$ at $B^4$ and
$13/65536$ at each of $A^2D^2,A^2C^2$.  Their distinct recipient
groups have original weights $1/256,9/64,3/8$, so all additions fit.

This proves every tensor comparison.  For positive $A,B,C,D$, it is
strict: the group containing $A^2B^2$ has a positive maximum and final
weight below one.  Its original weight is $3/8$; the outer-pivot transfers
leave it at $217/512$ or $97/256$, the first three inner-pivot transfers
leave it at $1/2$, and the inner ratio $4/3$ leaves it unchanged.
Together with the twelve three-corner profiles, this proves
Proposition~\ref{c-parallelogram}.  The finite maximum in
\eqref{c-parallelogram-norm} preserves strictness.

\section{Four-term arithmetic progressions}
\label{sec:ap4}
Let \(F=\{0,1,2,3\}\). For \(k=(k_0,k_1,k_2,k_3)\ge0\), write

\[
T(k)=T_F(k)=\sum_{s=0}^{12}
\max_{i_1+i_2+i_3+i_4=s}\prod_{r=1}^{4}k_{i_r},
\qquad i_r\in F.
\]

For a tuple \(v\), let \(v^R_i=v_{3-i}\). Define the five squared profiles

\[
\begin{aligned}
A&=(1,1/4,0,1/4),\\
B&=(1,0,1/4,1/4),\\
C&=(1,1/4,9/64,25/256),\\
D&=(1/4,1,0,1/4),\\
E&=(0,1,1/4,9/64),
\end{aligned}
\]

and the ten-element set

\[
\mathcal V=\{A,B,C,D,E,A^R,B^R,C^R,D^R,E^R\}.
\]

\begin{proposition}\label{prop:ap4}
The weighted norm has the finite-profile formula

\begin{equation}\label{eq:a1}
\boxed{\quad
\sup_{x\ge0,\;x*x\le1}\sum_{i=0}^{3}k_ix_i^2
=\max_{v\in\mathcal V}k\cdot v,
\quad}
\end{equation}

where the convolution inequality is coefficientwise at all seven indices, and

\begin{equation}\label{eq:a2}
\boxed{\quad
(k\cdot v)^4\le T(k)\qquad(v\in\mathcal V).
\quad}
\end{equation}

Consequently,

\begin{equation}\label{eq:a3}
x*x\le1\quad\Longrightarrow\quad
\left(\sum_i k_ix_i^2\right)^4\le T(k)
\qquad(k\ge0).
\end{equation}

For positive \(k_i\), comparison \eqref{eq:a2}, and hence \eqref{eq:a3} at the maximizing value,
is strict. These statements hold on every four-term integer arithmetic
progression after relabeling the coordinates.
\end{proposition}

\subsection{Feasibility and a six-profile convex hull}

Each member of \(\mathcal V\) is the square of a feasible tuple. For the five
unreversed members, the seven coefficients of \(\sqrt v*\sqrt v\) are

\begin{center}
\begin{tabular}{ll}
\toprule
Profile & Ordinary convolution coefficients\\
\midrule
\(A\) & \((1,1,1/4,1,1/2,0,1/4)\)\\
\(B\) & \((1,0,1,1,1/4,1/2,1/4)\)\\
\(C\) & \((1,1,1,1,29/64,15/64,25/256)\)\\
\(D\) & \((1/4,1,1,1/2,1,0,1/4)\)\\
\(E\) & \((0,0,1,1,1,3/8,9/64)\)\\
\bottomrule
\end{tabular}
\end{center}

Reversal reverses this coefficient list. Thus the right-hand side of \eqref{eq:a1} is
always attained by a feasible squared profile.

Let

\[
\mathcal S=\{A,B,D,A^R,B^R,D^R\}.
\]

Its convex hull has the exact elementary description

\begin{equation}\label{eq:a4}
\operatorname{conv}(\mathcal S)=
\left\{y\ge0:
\sum_i y_i=\frac32,\quad y_0,y_3\ge\frac14,\quad
y_0+y_3\le\frac54\right\}.
\end{equation}

To check the reverse inclusion, put

\[
r_0=\frac{4y_0-1}{3},\qquad
r_3=\frac{4y_3-1}{3},\qquad r_m=1-r_0-r_3.
\]

These are nonnegative and sum to one. Allocate the same fraction
\(p=y_1/(y_1+y_2)\) to the first member of each pair
\((A,B)\), \((B^R,A^R)\), and \((D,D^R)\), with pair masses
\(r_0,r_3,r_m\), respectively. The resulting convex combination is \(y\).
The denominator is positive because \(y_1+y_2\ge1/4\).

A useful consequence is the following sufficient domination test. If a tuple
\(b\ge0\) satisfies

\[
b_0,b_3\le\frac12,\qquad m=b_1+b_2\le1,
\]

\begin{equation}\label{eq:a5}
m+\max(b_0,1/4)+\max(b_3,1/4)\le\frac32,
\end{equation}

then some \(y\in\operatorname{conv}(\mathcal S)\) dominates \(b\)
coordinatewise. Indeed choose the final middle sum to be
\(m'=\max(m,1/4)\), and the endpoint sum to be \(3/2-m'\). The latter lies
in \([1/2,5/4]\) and is at least
\(\max(b_0,1/4)+\max(b_3,1/4)\): this is \eqref{eq:a5} when \(m\ge1/4\), and
follows from \(b_0,b_3\le1/2\) otherwise. Increase coordinates to these sums
and apply \eqref{eq:a4}.

\subsubsection{Self-contained proof of the sharp norm lemma}

Let \(N=\sum_i x_i^2\) and \(M=\max_i x_i^2\). The constraints imply
\(x_i^2\le1\), \(2x_ix_j\le1\) for \(i\ne j\), and
\(x_0x_3+x_1x_2\le1/2\). We repeatedly use

\begin{equation}\label{eq:a6a}
0\le r,s\le J\quad\Longrightarrow\quad r+s\le J+rs/J,
\end{equation}

which follows from \((J-r)(J-s)\ge0\).

If \(M\le1/2\), put \(r=x_0x_3\), \(s=x_1x_2\). Applying \eqref{eq:a6a} to the
squared coordinates in each pair yields

\[
N\le1+2r^2+2s^2\le1+2(r+s)^2\le3/2.
\]

Suppose \(1/2\le M\le3/4\), attained at index \(i\). Put
\(y=x_ix_{3-i}\), \(J=1/(4M)\), and let \(r\) be the product of the two
coordinates outside this pair. The other three squared coordinates are at most
\(J\), while \(r+y\le1/2\) and \(r\le J\). Hence

\[
N\le M+y^2/M+J+4M\min(J,1/2-y)^2.
\]

For \(0\le y\le1/2-J\), this expression increases with \(y\). For
\(1/2-J\le y\le1/2\), it is a convex quadratic. Its maximum is therefore
the larger of the two endpoint values

\[
B(M)=M+\frac1{2M}+\frac{(2M-1)^2}{16M^3},\qquad M+\frac1{2M},
\]

namely \(B(M)\). With \(w=2M-1\in[0,1/2]\),

\[
16M^3(3/2-B(M))=w(1-w^2-w^3)\ge0.
\]

Suppose \(M\ge3/4\) is at index one; index two follows by reversal. Put
\(J=1/(4M)\). The constraint \(M+2x_0x_2\le1\) and \eqref{eq:a6a} imply

\[
x_0^2+x_2^2\le J+M(1-M)^2,
\qquad N\le M+\frac1{2M}+M(1-M)^2.
\]

The difference from \(3/2\) is

\[
\frac{(1-M)(2M^3-2M^2+2M-1)}{2M}\ge0.
\]

The cubic increases, since its derivative is
\(6(M-1/3)^2+4/3>0\), and is \(7/32\) at \(M=3/4\).

Finally suppose an endpoint has squared value \(u\in[1/2,1]\). By reversal
take \(x_0^2=u\), and put
\(a=2x_0x_1\), \(b=2x_0x_2\), \(c=2x_0x_3\). The first four constraints
give

\[
0\le a\le1,\qquad 0\le b\le1-a^2/4,\qquad 0\le c\le1-ab/2.
\]

For fixed \(a\), the upper bound \(a^2+b^2+(1-ab/2)^2\) on
\(a^2+b^2+c^2\) is convex in \(b\). At \(b=0\), it is at most two. At
\(b=1-a^2/4\), it is

\[
\Psi(a)=2-a+3a^2/4+a^3/4-a^4/16+a^6/64.
\]

This is convex on \([0,1]\), since
\(\Psi''(a)=3/2+3a/2-3a^2/4+15a^4/32\ge3/4\), and its endpoint values
are \(2\) and \(125/64<2\). Therefore

\[
N=u+\frac{a^2+b^2+c^2}{4u}\le u+\frac1{2u}\le\frac32.
\]

These cases cover every feasible tuple. The norm bound is sharp at
\(x=(1/2,1,0,1/2)=\sqrt D\), as its feasibility table shows.

\subsection{Reduction to the ten profiles}

For a feasible \(x\), write \(b_i=x_i^2\). The convolution constraints give

\begin{equation}\label{eq:a6}
b_i\le1,\qquad b_ib_j\le\frac14\quad(i\ne j).
\end{equation}

We also use the previously proved sharp norm lemma

\begin{equation}\label{eq:a7}
\sum_i b_i\le\frac32.
\end{equation}

\subsubsection{All squared coordinates at most one half}

Suppose \(b_i\le1/2\) for all \(i\). Each consecutive triple has sum at most
\(9/8\). For example, the constraint
\(b_1+2\sqrt{b_0b_2}\le1\), together with \(b_0,b_2\le1/2\), gives

\[
b_0+b_2\le\frac12+2b_0b_2
\le\frac12+\frac{(1-b_1)^2}{2},
\]

\begin{equation}\label{eq:a8}
b_0+b_1+b_2\le1+\frac{b_1^2}{2}\le\frac98.
\end{equation}

The other triple is the reversed argument.

Condition \eqref{eq:a5} now holds. If both endpoints exceed \(1/4\), use \eqref{eq:a7}; if both
are below \(1/4\), use \(b_1+b_2\le1\). If just one is below \(1/4\),
the expression in \eqref{eq:a5} is a consecutive triple sum plus \(1/4\), at most
\(11/8<3/2\). Thus \(b\) is dominated by the convex hull of \(\mathcal S\).

\subsubsection{A middle squared coordinate at least one half}

By reversal, it suffices to write

\[
b=(a,u,v,d),\qquad 1/2\le u\le1.
\]

Put \(J=1/(4u)\). By \eqref{eq:a6}, \(a,v,d\le J\le1/2\). The constraint at
index two gives \(\sqrt{av}\le(1-u)/2\). Consequently

\[
a+v\le J+\frac{av}{J}\le\frac1{4u}+u(1-u)^2.
\]

The following bound is useful throughout this case:

\begin{equation}\label{eq:a9}
a+u+v\le\frac54.
\end{equation}

Indeed
\[
\frac54-\left(u+\frac1{4u}+u(1-u)^2\right)
=\frac{(1-u)(4u^3-4u^2+4u-1)}{4u}\ge0.
\]

The cubic is increasing, because \(12u^2-8u+4>0\), and its value at
\(u=1/2\) is \(1/2\).

First suppose \(m=u+v\le1\). The domination test \eqref{eq:a5} holds. The cases when
both endpoints are on the same side of \(1/4\) follow from \eqref{eq:a7} or \(m\le1\).
If \(a\ge1/4\) and \(d\le1/4\), use \eqref{eq:a9}. If \(a\le1/4\) and
\(d\ge1/4\), the index-four constraint gives

\[
u+v+d\le1+(\sqrt u-\sqrt d)^2\le\frac54,
\]

because \(d\le1/2\le u\le1\) and \(d\ge1/4\), so
\(0\le\sqrt u-\sqrt d\le1/2\). Thus \eqref{eq:a5} again proves domination by
\(\operatorname{conv}(\mathcal S)\).

Now suppose \(m=u+v>1\). Define

\[
\theta=4(m-1),\qquad \beta=4-3u-4v,\qquad \gamma=1-u.
\]

All three numbers are nonnegative and their sum is one. For \(\beta\), use
\(v\le1/(4u)\) to obtain

\[
\beta\ge4-3u-1/u=\frac{(1-u)(3u-1)}{u}\ge0.
\]

The convex combination

\begin{equation}\label{eq:a10}
y=\theta E+\beta D+\gamma D^R
=\left(\frac54-m,\ u,\ v,\ \frac{11-7m}{16}\right)
\end{equation}

dominates \(b\). The first coordinate does so by \eqref{eq:a9}. For the last, the
index-four constraint gives \(d\le(1-v)^2/(4u)\). We claim

\begin{equation}\label{eq:a11}
\frac{(1-v)^2}{4u}\le\frac{11-7u-7v}{16}
\qquad\left(1-u\le v\le\frac1{4u}\right).
\end{equation}

After multiplying by \(16u\), the required nonnegative expression is

\[
G_u(v)=11u-7u^2-7uv-4+8v-4v^2.
\]

It is concave in \(v\), so its minimum on this interval is at an endpoint.
The two endpoint values are

\[
G_u(1-u)=4u(1-u),
\]

\[
G_u(1/(4u))
=\frac{(1-u)(28u^3-16u^2+7u-1)}{4u^2}.
\]

The latter cubic is increasing on \([1/2,1]\), since its derivative
\(84u^2-32u+7\) is positive there, and its value at \(1/2\) is two.
This proves \eqref{eq:a11}, then \eqref{eq:a10}, and completes the middle-coordinate case.

\subsubsection{An endpoint squared coordinate at least one half}

By reversal, assume \(u=x_0^2\in[1/2,1]\). Put

\[
\alpha=2x_0x_1,\qquad \beta=2x_0x_2,\qquad \gamma=2x_0x_3.
\]

The first four convolution constraints give

\[
0\le\alpha\le1,\qquad
0\le\beta\le1-\alpha^2/4,\qquad
0\le\gamma\le1-\alpha\beta/2.
\]

For arbitrary nonnegative \(k\),

\[
k\cdot b
=k_0u+\frac{k_1\alpha^2+k_2\beta^2+k_3\gamma^2}{4u}.
\]

For fixed \(\alpha\), the numerator is at most

\[
k_1\alpha^2+k_2\beta^2+k_3(1-\alpha\beta/2)^2,
\]

which is convex in \(\beta\). At \(\beta=0\), it is at most \(k_1+k_3\).
At \(\beta=1-\alpha^2/4\), it is the polynomial

\[
\begin{aligned}
\Phi(\alpha)={}&k_2+k_3-k_3\alpha
+(k_1-k_2/2+k_3/4)\alpha^2
+k_3\alpha^3/4\\
&+(k_2/16-k_3/8)\alpha^4+k_3\alpha^6/64.
\end{aligned}
\]

This polynomial has its maximum on \([0,1]\) at an endpoint, even when it is
not convex. To see this, observe

\[
\Phi'(0)=-k_3\le0,
\]

\[
\Phi'''(\alpha)=\frac32 k_2\alpha
+\frac38 k_3(4-8\alpha+5\alpha^3)\ge0.
\]

For the last inequality, the minimum of the cubic on \([0,1]\) is at
\(\alpha=\sqrt{8/15}\), where its value is
\(4-(16/3)\sqrt{8/15}\ge0\), since \(\sqrt{8/15}\le3/4\).
Thus \(\Phi'\) is convex and starts nonpositive; its nonpositive sublevel
set is an interval containing zero. Accordingly \(\Phi\) first decreases
and then increases, with either portion possibly absent. Its maximum is at an
endpoint.

The two endpoint values of \(\Phi\) are \(k_2+k_3\) and
\(k_1+9k_2/16+25k_3/64\). Therefore the numerator is at most

\[
H=\max\{k_1+k_3,\ k_2+k_3,\ k_1+9k_2/16+25k_3/64\}.
\]

For each of these three linear forms, convexity in \(u\in[1/2,1]\) reduces
\(k_0u+H/(4u)\) to its two endpoint values. At \(u=1\), the resulting
squared profiles are \(A,B,C\). At \(u=1/2\), they are

\[
W_A=(1/2,1/2,0,1/2),\qquad
W_B=(1/2,0,1/2,1/2),
\]

\[
W_C=(1/2,1/2,9/32,25/128).
\]

The first two belong to the sparse convex hull, since

\[
W_A=(A+D+B^R)/3,\qquad
W_B=(B+D^R+A^R)/3.
\]

The last is dominated coordinatewise by

\[
\frac13 C+\frac5{12}D+\frac14D^R
=\left(\frac12,\frac12,\frac{19}{64},\frac{51}{256}\right)
\ge W_C.
\]

Consequently \(k\cdot b\le\max_{v\in\mathcal V}k\cdot v\) in this case
as well. Every feasible tuple falls into one of the three cases. Together with
the explicit feasibility table, this proves \eqref{eq:a1}.

\subsection{Comparing the ten linear forms with the four-fold maximum}

Reversal preserves \(T\), so it suffices to handle \(A,B,C,D,E\).

\subsubsection{A coefficient lemma; profiles C and E}

If \(v\ge0\) has the property that every coefficient of
\((\sum_i v_i z^i)^4\) is at most one, then

\begin{equation}\label{eq:a12}
(k\cdot v)^4
=\sum_s\sum_{i_1+\cdots+i_4=s}\prod_r(v_{i_r}k_{i_r})
\le\sum_s\max_{i_1+\cdots+i_4=s}\prod_r k_{i_r}=T(k).
\end{equation}

The nonnegative formal series

\[
(1-z)^{-1/4}=1+z/4+5z^2/32+15z^3/128+\cdots
\]

coefficientwise dominates
\(1+z/4+9z^2/64+25z^3/256\). Raising to the fourth power preserves this
order, and \(((1-z)^{-1/4})^4=(1-z)^{-1}\) has all coefficients equal to
one. This proves the required coefficient property for \(C\).

For \(E\), use
\(\sum_i E_i z^i=z(1+z/4+9z^2/64)\). The same comparison applies before
the harmless fourth-power shift by four. Thus \eqref{eq:a12} proves both profiles.

\subsubsection{Profile B}

Set \(a=k_0\), \(b=k_2\), \(c=k_3\), and ignore \(k_1\); doing so can
only decrease \(T\). The target linear form is \(a+b/4+c/4\), on indices
\(0,2,3\). Grouping its fourth-power expansion by index sum, the only fibers
with more than one monomial are

\begin{center}
\begin{tabular}{lll}
\toprule
Index sum & Monomials and their expansion weights & Total weight\\
\midrule
6 & \(ab^3/16+3a^2c^2/8\) & \(7/16\)\\
8 & \(b^4/256+3abc^2/16\) & \(49/256\)\\
9 & \(b^3c/64+ac^3/16\) & \(5/64\)\\
\bottomrule
\end{tabular}
\end{center}

Every other fiber has one monomial with weight at most one. Each fiber is thus
bounded by its largest kernel monomial. Summing proves the comparison for \(B\).

\subsubsection{The common three-point maximum for A and D}

Now put \(a=k_0\), \(b=k_1\), \(c=k_3\), and ignore \(k_2\). On the
indices \(\{0,1,3\}\), the four-fold maximum sum is exactly

\begin{equation}\label{eq:a13}
\begin{aligned}
\mathcal T(a,b,c)={}&a^4+a^3b+a^2b^2
+\max(a^3c,ab^3)+\max(a^2bc,b^4)+ab^2c\\
&+\max(a^2c^2,b^3c)+abc^2+b^2c^2+ac^3+bc^3+c^4.
\end{aligned}
\end{equation}

These terms correspond to sums \(0,1,2,3,4,5,6,7,8,9,10,12\), respectively;
sum eleven is absent. In particular \(T(k)\ge\mathcal T(a,b,c)\).

For \(A\), let
\[
\begin{aligned}
P_A={}&a^4+a^3b+a^3c+a^2b^2
+\tfrac34a^2bc+\tfrac14b^4+ab^2c\\
&+\tfrac12a^2c^2+\tfrac12b^3c
+abc^2+b^2c^2+ac^3+bc^3+c^4.
\end{aligned}
\]

Selecting one term or a convex average in each maximum in \eqref{eq:a13} gives
\(\mathcal T\ge P_A\). An exact nonnegative remainder is

\begin{equation}\label{eq:a14}
\begin{aligned}
P_A-(a+b/4+c/4)^4={}&\frac{b^2(a-b)^2}{32}
+\frac{19a^2b^2}{32}+\frac{a^2c^2}{8}\\
&+\frac{13ab^2c}{16}+\frac{13abc^2}{16}
+\frac{15ac^3}{16}+\frac{55b^4}{256}\\
&+\frac{31b^3c}{64}+\frac{125b^2c^2}{128}
+\frac{63bc^3}{64}+\frac{255c^4}{256}.
\end{aligned}
\end{equation}

Thus \eqref{eq:a2} holds for \(A\).

For \(D\), choose the terms \(ab^3,b^4,b^3c\) in the three maxima of
\eqref{eq:a13}, obtaining the lower bound

\[
P_D=a^4+a^3b+a^2b^2+ab^3+b^4+ab^2c+b^3c
+abc^2+b^2c^2+ac^3+bc^3+c^4.
\]

In comparison with \((a/4+b+c/4)^4\), the only missing monomials are

\[
\frac{a^3c}{64},\qquad\frac{3a^2bc}{16},\qquad\frac{3a^2c^2}{128}.
\]

The arithmetic--geometric mean inequalities

\[
a^3c\le\frac{3a^4+c^4}{4},\qquad
a^2bc\le\frac{a^4+b^2c^2}{2},\qquad
a^2c^2\le\frac{a^4+c^4}{2}
\]

absorb their total using at most

\[
\frac{30}{256}a^4+\frac4{256}c^4+\frac3{32}b^2c^2.
\]

The available coefficient surpluses of these three monomials in
\(P_D-(a/4+b+c/4)^4\) are respectively
\(255/256,255/256,5/8\), all larger than the required amounts. All other
coefficient surpluses are nonnegative. This proves the comparison for \(D\),
and completes \eqref{eq:a2}.

For positive \(k_i\), these comparisons are strict. The profile \(C\) has a
strict coefficient deficit at index sum two. Profiles \(A,D\) omit the positive
sum-eleven contribution involving \(k_2\); \(B\) omits the positive sum-one
contribution involving \(k_1\); and \(E\) omits the positive sum-zero
contribution. Reversal handles the other five profiles.

\section{One three-term progression}
\label{sec:isolated}
Let

\[
F=\{0,a,2a,b\},\qquad a>0,
\]

where the four integers are distinct and the only nontrivial equality of pair
sums is \(0+2a=a+a\). Equivalently, \(q=b/a\) is not any of

\begin{equation}\label{eq:i1}
-2,-1,0,1/2,1,3/2,2,3,4.
\end{equation}

Coordinates are \((L,M,R,O)=(0,a,2a,b)\).

Let \(\mathcal V\) consist of the following eight squared profiles:

\begin{equation}\label{eq:i2}
\begin{array}{ll}
V_L=(1,1/4,9/64,1/4),&V_R=(9/64,1/4,1,1/4),\\
D_L=(1/4,1,0,1/4),&D_R=(0,1,1/4,1/4),\\
E_L=(1,0,1/4,1/4),&E_R=(1/4,0,1,1/4),\\
O_1=(1/4,1/4,1/4,1),&O_3=(1/3,1/3,1/3,3/4).
\end{array}
\end{equation}

\begin{proposition}\label{prop:isolated}
For arbitrary \(k\ge0\), the weighted norm is

\begin{equation}\label{eq:i3}
\sup_{x*x\le1}\sum_{p\in F}k(p)x(p)^2
=\max_{v\in\mathcal V}k\cdot v.
\end{equation}

In particular the sharp unweighted norm is \(7/4\), attained by \(O_1,O_3\).
Each listed profile satisfies

\begin{equation}\label{eq:i4}
(k\cdot v)^4\le T_F(k)\quad\text{for all eight profiles},
\end{equation}

strictly when all four kernel values are positive.
\end{proposition}

\subsection{Exact weighted norm reduction}

Write \(w_p=x(p)^2\). Ordinary feasibility on this support is exactly the
individual bounds \(w_p\le1\), pair bounds \(w_pw_q\le1/4\), and

\begin{equation}\label{eq:i5}
w_M+2\sqrt{w_Lw_R}\le1.
\end{equation}

The pair bounds within the collided fiber are redundant but valid. Every
profile in \eqref{eq:i2} is feasible, by direct substitution in \eqref{eq:i5} and the pair bounds.

We use the following elementary maximization repeatedly. For
\(0\le l,r\le J\), \(lr\le c\le J^2\), the maximum of any nonnegative
linear objective occurs at \((J,c/J)\) or \((c/J,J)\). On the part of the
boundary with \(r=c/l\), the objective is convex in \(l\); elsewhere one
coordinate can be increased. When \(c\ge J^2\), use \((J,J)\).

First suppose every \(w_p\le1/2\). Set \(w_O=1/2\). At fixed
\(m=w_M\), the endpoint rule gives an endpoint \(1/2\), the other
\((1-m)^2/2\). The objective is convex in \(m\in[0,1/2]\), so this region
reduces to

\begin{equation}\label{eq:i6}
H_0=(1/2,0,1/2,1/2),\quad H_L=(1/2,1/2,1/8,1/2),\quad H_R,
\end{equation}

where \(H_R\) exchanges the endpoints of \(H_L\).

Otherwise let \(u\in[1/2,1]\) be a largest coordinate and put \(J=1/(4u)\).
All other coordinates are at most \(J\), so all their mutual pair bounds
are automatic.

If the largest coordinate is \(L\), set \(O=J\). For fixed \(M=m\),
the best right endpoint is \(J(1-m)^2\). Convexity in \(m\in[0,J]\)
reduces to

\begin{equation}\label{eq:i7}
(u,0,J,J),\qquad (u,J,J(1-J)^2,J).
\end{equation}

The first objective is convex in \(u\), giving \(H_0,E_L\). For the second,
set

\[
C=(1/u-1)/3,\quad B=\frac{64}{55}(3/2-u-1/(2u)),\quad A=1-2C-B.
\]

These coefficients are nonnegative:
\(B=32(2u-1)(1-u)/(55u)\), and
\(A=(192u^2-13u-14)/(165u)>0\) on this interval. The convex combination

\[
A V_L+B V_R+C D_L+C O_1
\]

matches the \(L,M,O\) coordinates of the second profile in \eqref{eq:i7}. Its right
coordinate is \(317/192-u-49/(96u)\). The excess over \(J(1-J)^2\),
multiplied by \(192u^3\), is

\[
(1-u)(192u^3-125u^2+21u-3)\ge0.
\]

The cubic is \(1/4\) at \(u=1/2\) and increases: its derivative is
\(576u^2-250u+21\), positive there and increasing thereafter. The case of
largest coordinate \(R\) is reflected. This also dominates \(H_L,H_R\)
by setting \(u=1/2\).

If \(M=u\) is largest, set \(O=J\). The endpoint rule gives
\((L,R)=(J,u(1-u)^2)\), or its reflection. Put the larger endpoint kernel
weight on \(J\), say \(k_L\ge k_R\). The objective is

\[
k_Mu+k_R[J+u(1-u)^2]+(k_L-k_R+k_O)J.
\]

It is convex in \(u\). Besides \(J''>0\), the needed derivative is
\([J+u(1-u)^2]''=1/(2u^3)+6u-4>0\): split at \(u=2/3\), and use
\(1/(2u^3)\ge27/16\), \(6u-4\ge-1\) on the lower interval.
The endpoints give \(H_L,D_L\), or their reflections.

If \(O=u\) is largest and \(u\ge3/4\), the best remaining coordinates are
all \(J\); the objective is convex in \(u\), giving \(O_3,O_1\). For
\(1/2\le u\le3/4\), optimize at fixed midpoint and use convexity in that
midpoint. The only profiles needed are

\begin{equation}\label{eq:i8}
(J,1-2J,J,u),\qquad (J,J,u-1/2+J/4,u),
\end{equation}

and the reflected second profile. The first is dominated by a convex
combination of \(H_0,O_3\): use weight \(3(1-2J)\) on \(O_3\).
This matches the first three coordinates and gives outsider coordinate
\(5/4-3J/2\ge u\), since \((u-1/2)(u-3/4)\le0\). The second objective
in \eqref{eq:i8} is a nonnegative multiple of \(1/u\) plus a linear function, hence
convex, giving \(H_L,O_3\).

Finally \(H_0\le(E_L+E_R+O_1)/3\) coordinatewise. This completes the
upper bound in \eqref{eq:i3}; feasibility proves equality. The profile sums,
in the order listed, are
\[
105/64,105/64,3/2,3/2,3/2,3/2,7/4,7/4.
\]

\subsection{The four triangle profiles}\label{sec:i-triangles}

Each support triangle of \(D_L,D_R,E_L,E_R\) is 2-sum-Sidon: an additional
midpoint relation would violate the isolated-AP3 assumption. For a nonnegative
kernel on any 2-sum-Sidon triangle, let \(a\) be its largest value and \(R\)
the sum of the other two. Fixing two copies of the largest point in the
four-fold maximum gives

\[
T(k)\ge a^4+a^3R+a^2h_2(\text{other two values})
\ge a^4+a^3R+\frac23a^2R^2.
\]

Here \(h_2\) is the sum of the squares and pair product. Since \(R\le3a\),

\[
(a+R/4)^4\le a^4+a^3R+\frac{153}{256}a^2R^2.
\]

The difference is at least \(53a^2R^2/768\), positive when two values are
positive. The largest pivot maximizes the unit-and-two-quarter linear form.
Restricting \(k\) to the triangle and using monotonicity proves strict \eqref{eq:i4}
for these four profiles when the full kernel is positive.

\subsection{Common quartic framework and collision inventory}

Write kernel values as \((A,B,C,D)=(k_L,k_M,k_R,k_O)\). In the fourth
power of a weighted linear form, use triangular grid coordinates

\[
i=n_B+2n_C,\qquad j=n_D,\qquad
0\le j\le4,\quad0\le i\le2(4-j).
\]

All monomials in this grid point have actual sum \(ia+jb\). If the total
nonnegative coefficient weights of a projected fiber are at most one, its
weighted polynomial is bounded by that fiber's kernel-product maximum.
Summing gives the required quartic comparison. The argument also allows
stated AM--GM redistributions to other fibers before applying this test.

For \(O_1,O_3\), reflection about the AP midpoint permits \(q=b/a>1\).
Write \(q=r/s\) in lowest terms. Additional grid collisions differ by
\((r,-s)\), so require \(r\le8,s\le4\). Removing the forbidden ratios
in \eqref{eq:i1} leaves precisely

\begin{equation}\label{eq:i9}
\begin{array}{c|l}
s& r\\\hline
1&5,6,7,8\\
2&5,7\\
3&4,5,7,8\\
4&5,7.
\end{array}
\end{equation}

Each projected fiber contains at most two grid points. Ratios outside this
list have no additional fourth-sum collision.

For the endpoint profile \(V_L\), also consider its reflection \(V_R\) by
allowing \(q<1\) in the same template. The additional positive ratios are
\(1/3,2/3,1/4,3/4\). For \(q=-r/s<0\), a collision increases both grid
coordinates by \((r,s)\), so \(r+2s\le8\). The allowed negative ratios are

\begin{equation}\label{eq:i10}
-3,-4,-5,-6,-1/2,-3/2,-1/3,-2/3.
\end{equation}

Again no fiber has more than two grid points. These lists exhaust the
reflected cases as well: reflection replaces \(q\) by \(2-q\).

We repeatedly use

\begin{equation}\label{eq:i11}
U+\epsilon V\le\max(U,V)+\epsilon\min(U,V)\quad(0\le\epsilon\le1)
\end{equation}

and

\begin{equation}\label{eq:i12}
\min(X^3Y,Z^4)\le X^2Z^2+X^2Y^2.
\end{equation}

For \eqref{eq:i12}, divide by \(X^4\) when \(X>0\). The assertion
\(\min(y,z^4)\le z^2+y^2\) follows by splitting into \(z\le1\),
\(z\ge1,y\le1\), and \(y\ge1\). The zero case is immediate.

\subsection{The profile O3}\label{sec:i-o3}

Here \(E=(A+B+C)/3+3D/4\). Let

\[
P_2=(1,2,3,2,1),\quad P_3=(1,3,6,7,6,3,1),
\]

\[
P_4=(1,4,10,16,19,16,10,4,1)
\]

be the coefficients of \((1+z+z^2)^2,(1+z+z^2)^3,(1+z+z^2)^4\).
The grid coefficient rows of \(E^4\) are

\[
\begin{gathered}
j=0:P_4/81;\quad j=1:P_3/9;\quad j=2:3P_2/8;\\
j=3:(9/16,9/16,9/16);\quad j=4:(81/256).
\end{gathered}
\]

At grid \((2,2)\), move weight \(1/4\) of \(ACD^2\) using
\(ACD^2\le(A^2D^2+C^2D^2)/2\). The new \(j=2\) row is

\[
(1/2,3/4,7/8,3/4,1/2).
\]

All other rows are unchanged. Uncollided fibers have weight at most \(7/8\).
For the four denominator rows of \eqref{eq:i9}, upper bounds on collided-fiber weights
are respectively

\[
31/36,\quad 599/648,\quad755/768,\quad81/256+16/81,
\]

all below one. These follow by adding the displayed rows with displacement
\((r,-s)\). For example the closest case is \(r/s=4/3\), where grid
\((0,4)\) meets \((4,1)\), giving \(81/256+2/3=755/768\).
The coefficient test proves \eqref{eq:i4}, with strictness for every nonzero kernel.

\subsection{The profile O1}\label{sec:i-o1}

Now \(E=D+(A+B+C)/4\). Its coefficient rows are

\[
j=0:P_4/256;\quad j=1:P_3/16;\quad j=2:3P_2/8;
\quad j=3:(1,1,1);\quad j=4:(1).
\]

Move weight \(1/2\) of \(ACD^2\) by the same AM--GM inequality. The \(j=2\)
row becomes

\begin{equation}\label{eq:i13}
(5/8,3/4,5/8,3/4,5/8).
\end{equation}

For ratios \(5,6,7,8,7/2\), every projected fiber now has weight at most
one. At the other ratios, retain the full \(j=3\) or \(j=4\) polynomial
in each overloaded fiber and remove its other grid-point polynomial. The
following ledgers absorb precisely those removed terms. All unmentioned
fibers pass the coefficient test with \eqref{eq:i13}.

At \(q=5/2\), the removed terms are

\[
3DBC^2/16+DC^3/16.
\]

They are bounded by \(D^2C^2/8+3B^2C^2/32+C^4/32\). The original
weights of these three projected fibers were \(5/8,101/128,193/256\).
After allocation they are \(3/4,113/128,201/256\), below one.

At \(q=5/3\), the removed AP-only rows are at \(i=5,6,7\), with total
coefficient weight \(30/256\); the removed \(j=1,i=5\) term is
\(3DBC^2/16\). Every removed unweighted monomial is at most

\[
S=A^4+B^4+C^4+D^2C^2,
\]

by weighted AM--GM (for \(DBC^2\), use
\(DBC^2\le D^2C^2/2+B^4/4+C^4/4\)). The total removed weight is
\(39/128\). The four \(S\) monomials occupy distinct fibers, of original
weights \(1/256,19/256,1/256,5/8\). Adding \(39/128\) to each remains
below one; the largest becomes \(119/128\).

At \(q=4/3\), the removed AP-only rows are \(i=4,5,6\), with total
weight \(45/256\). Bound them by
\(45(A^4+B^4+C^4)/256\). The removed \(j=1,i=4\) polynomial is
\(3DAC^2/16+3DB^2C/16\), bounded by

\[
3D^2C^2/16+3A^4/64+3C^4/64+3B^4/32.
\]

Thus the allocations are \(57/256\) each to \(A^4,C^4\), \(3/16\) to
\(D^2C^2\), and \(69/256\) to \(B^4\). The last shares the fiber of
\(D^3A\). Apply \eqref{eq:i11}-\eqref{eq:i12} to
\(D^3A+(69/256)B^4\); the remainder is at most
\(69(D^2B^2+D^2A^2)/256\). The final affected weights are
\(58/256\) at each of \(A^4,C^4\), \(13/16\) at \(D^2C^2\), and
\(229/256\) at each of \(D^2B^2,D^2A^2\). All are below one.

At \(q=7/3\), the removed terms are \(BC^3/64+C^4/256\), at the fibers
of \(D^3A,D^3B\), respectively. They are bounded by
\(B^4/256+C^4/64\). Use \eqref{eq:i11}-\eqref{eq:i12} for \(D^3B+C^4/64\), allocating its
remainder to \((D^2C^2+D^2B^2)/64\). The \(B^4\) fiber had weight
\(19/256\); the two quadratic fibers had weight \(5/8\). These additions
fit. At \(q=8/3\), only \(C^4/256\) is removed, at the \(D^3A\) fiber;
\eqref{eq:i11}-\eqref{eq:i12} allocate its remainder to \((D^2C^2+D^2A^2)/256\).

At \(q=5/4\), the removed AP-only row is \(i=5\), of total weight
\(1/16\), at the \(D^4\) fiber. Each monomial is at most
\(A^4+B^4+C^4\). These three distinct unoverloaded fibers have weights
\(1/256,19/256,1/256\); adding \(1/16\) fits. At \(q=7/4\), the sole
removed term is \(BC^3/64\), absorbed into \((B^4+C^4)/64\) in the same
way. This exhausts \eqref{eq:i9}. Every ledger leaves positive slack in a quadratic
fiber, so \eqref{eq:i4} is strict for positive kernels.

\subsection{Endpoint profiles: coefficients and positive ratios}

It suffices to study \(V_L\) at all ratios in \eqref{eq:i9}-\eqref{eq:i10} and the four positive
ratios below one. Reflection then gives \(V_R\). Write

\[
E=A+B/4+9C/64+D/4.
\]

Its triangular coefficient rows, in increasing \(i\), are

\begin{equation}\label{eq:i14}
\begin{array}{c|l}
j&\text{row}\\\hline
0&(1,1,15/16,31/64,467/2048,279/4096,1215/65536,
729/262144,6561/16777216)\\
1&(1,3/4,39/64,29/128,351/4096,243/16384,729/262144)\\
2&(3/8,3/16,33/256,27/1024,243/32768)\\
3&(1/16,1/64,9/1024)\\
4&(1/256).
\end{array}
\end{equation}

These are the coefficients of \((1+z/4+9z^2/64+w/4)^4\). Without extra
collisions they all pass the coefficient test. For \(q>1\) in \eqref{eq:i9}, the test
also passes directly unless \(q\) is an integer \(5,6,7,8\). At those
integers the only overloaded fiber is \(A^3D\), paired with AP grid point
\((q,0)\). Its additional polynomial is as follows:

\begin{equation}\label{eq:i15}
\begin{array}{c|l}
q&\text{additional polynomial}\\\hline
5&243ABC^2/4096+9B^3C/1024\\
6&729AC^3/65536+243B^2C^2/32768\\
7&729BC^3/262144\\
8&6561C^4/16777216.
\end{array}
\end{equation}

For \(q=5\), use
\(ABC^2\le(A^2B^2+C^4)/2\) and
\(B^3C\le(3B^4+C^4)/4\). The extra weight at \(A^2B^2\) is
\(243/8192<1/16\), fitting its original weight \(15/16\). The extra
weights at \(B^4,C^4\) are each below \(1/16\); their original fiber
weights are \(467/2048\) and \(29/128+6561/16777216\), respectively.

For \(q=6\), use \(AC^3\le(A^2C^2+C^4)/2\) and
\(B^2C^2\le(B^4+C^4)/2\). The first and third terms named in these
bounds, \(A^2C^2,B^4\), share AP fiber four. Its extra total weight, and
the extra weight at \(C^4\), are both \(1215/131072<1/64\). Their
original weights are \(467/2048\) and \(39/64+6561/16777216\).

For \(q=7\), use \(BC^3\le(B^4+3C^4)/4\); both extra allocations are
below \(1/256\), fitting original weights \(467/2048\) and
\(3/4+6561/16777216\). For \(q=8\), apply \eqref{eq:i11}-\eqref{eq:i12} directly to the
overloaded fiber \(A^3D+\epsilon C^4\), where
\(\epsilon=6561/16777216<1/256\), allocating the remainder to
\(\epsilon(A^2C^2+A^2D^2)\). Those fibers have original weights
\(467/2048\) and \(3/8\).

For the positive ratios below one, \(q=2/3,3/4\) pass directly using \eqref{eq:i14}.
At \(q=1/3\), the only overload is \(A^3B+AD^3/16\). Allocate the second
term to \((A^2D^2+D^4)/32\), in fibers originally of weights
\(3/8\) and \(3/4+1/256\). At \(q=1/4\), the sole overload is
\(A^3B+D^4/256\). Equations \eqref{eq:i11}-\eqref{eq:i12} allocate its remainder to
\((A^2D^2+A^2B^2)/256\), fitting original weights \(3/8,15/16\).

\subsection{Endpoint profiles: negative ratios}

The following table lists exactly the additional, high-grid-point polynomials
removed from overloaded fibers. Retain their low-grid-point polynomials in
full. All other fibers already pass the coefficient test with \eqref{eq:i14}.

\begin{center}
\begin{tabular}{ll}
\toprule
q & Removed polynomial\\
\midrule
-3 & \(B^3D/64+27ABCD/128+27B^2CD/1024+243AC^2D/4096+27BCD^2/1024\)\\
-4 & \(27B^2CD/1024+243AC^2D/4096+243BC^2D/16384+243C^2D^2/32768\)\\
-5 & \(243BC^2D/16384+729C^3D/262144\)\\
-6 & \(729C^3D/262144\)\\
-1/2 & \(3ABD^2/16+3B^2D^2/128+27ACD^2/256+BD^3/64\)\\
-3/2 & \(27BCD^2/1024+243C^2D^2/32768\)\\
-1/3 & \(BD^3/64+9CD^3/1024\)\\
-2/3 & \(9CD^3/1024\)\\
\bottomrule
\end{tabular}
\end{center}

For \(q=-3,-4\), every removed monomial is at most

\[
S=A^2C^2+B^4+C^4+D^4+B^2D^2+A^2D^2.
\]

Use weighted AM--GM for monomials without \(A\), and
\(ABCD\le(A^2C^2+B^2D^2)/2\),
\(AC^2D\le(A^2D^2+C^4)/2\). The removed total weights are respectively
\(1387/4096<11/32\) and \(3537/32768<1/8\). Allocate \(11/32\), or
\(1/8\), to each monomial of \(S\). The first two share a fiber, originally
of weight \(467/2048\); even the larger addition gives
\(467/2048+2(11/32)=1875/2048<1\). The other four occupy distinct
fibers, originally of weights
\(6561/16777216,1/256,33/256,3/8\), so all allocations fit.

For \(q=-1/2\), use

\[
S=A^2D^2+B^4+C^4+D^4.
\]

Every removed monomial is at most \(S\): for example
\(ABD^2\le A^2D^2/2+B^4/4+D^4/4\), and the same bound with \(C\)
replaces \(B\) for \(ACD^2\). The removed total weight is
\(85/256<3/8\). The four \(S\) monomials occupy distinct fibers whose
original weights are
\[
3/8,467/2048,6561/16777216,1/256.
\]
Adding \(3/8\) to each fits.

For each remaining row of the table, every monomial is at most
\(B^4+C^4+D^4\) by weighted AM--GM. The removed total weights have the
following convenient upper bounds:

\[
\begin{array}{c|ccccc}
q&-5&-6&-3/2&-1/3&-2/3\\\hline
\text{bound}&1/32&1/256&1/16&1/32&1/64.
\end{array}
\]

These three pure fourth powers occupy distinct unoverloaded fibers, of
original weights \(467/2048,6561/16777216,1/256\), so the stated additions
fit. This finishes \eqref{eq:i10} and hence both endpoint profiles.

The endpoint inequalities are strict for positive kernels. Except at
\(q=2/3\), the AP-sum-two fiber still has total coefficient weight below
one after the allocations. At \(q=2/3\), use the AP-sum-four fiber, whose
weight is \(467/2048+9/1024<1\). In each case its maximum is positive.
Together with Sections~\ref{sec:i-triangles}, \ref{sec:i-o3}, and
\ref{sec:i-o1}, this proves strict \eqref{eq:i4}.

\section{Two intersecting three-term progressions}\label{sec:double}

The support $F=\{0,1,2,4\}$ contains the progressions $\{0,1,2\}$ and
$\{0,2,4\}$. Both relations enter the weighted estimate. For a nonnegative
$x$ supported on $F$, put $w_i=x(i)^2$. The condition $x*x\leq 1$ is
equivalent to
\begin{equation}\label{d-feasibility}
\begin{gathered}
0\leq w_i\leq1,\qquad w_iw_j\leq\tfrac14\quad(i\ne j),\\
w_1+2\sqrt{w_0w_2}\leq1,\qquad
w_2+2\sqrt{w_0w_4}\leq1.
\end{gathered}
\end{equation}
All tuples in this section use the coordinate order $(0,1,2,4)$. We shall
bound nonnegative linear functionals of $w$ by their values on the profiles
\begin{equation}\label{d-profiles}
\begin{aligned}
H_0&=(\tfrac12,0,\tfrac12,\tfrac12),&
H_L&=(\tfrac12,\tfrac12,\tfrac18,\tfrac12),\\
H_R&=(\tfrac18,\tfrac12,\tfrac12,\tfrac12),&
V_L&=(1,\tfrac14,\tfrac9{64},\tfrac14),\\
V_R&=(\tfrac9{64},\tfrac14,1,\tfrac14),&
R&=(\tfrac9{64},\tfrac14,\tfrac14,1),
\end{aligned}
\end{equation}
and on feasible profiles with a zero coordinate. The tuples
$H_0,H_L,V_L,V_R$ are used only as upper bounds; they need not satisfy
\eqref{d-feasibility}.

\subsection{Profiles with a zero coordinate}

\begin{lemma}\label{d-boundary}
Suppose that $w$ satisfies \eqref{d-feasibility} and has a zero coordinate.
For every nonnegative kernel $k$ on $F$, put
$k_w=k\,1_{\operatorname{supp}(w)}$. Then
\begin{equation}\label{d-boundary-comparison}
(k\cdot w)^4\leq T_F(k_w)\leq T_F(k).
\end{equation}
The last comparison is strict when $k$ is positive on $F$. Moreover,
$\sum_iw_i\leq3/2$.
\end{lemma}

\begin{proof}
A three-term progression embeds in a four-term progression. Extending $k$
and $x$ by zero at the additional point preserves ordinary convolution
feasibility and the two quantities in the weighted comparison of
Section~\ref{sec:ap4}. This proves the first inequality when the remaining
support is a progression.

For a triangle with no three-term progression, its pair sums are distinct
apart from interchange. Write $a$ for its largest kernel value and $r$ for
the sum of the other two. Its weighted squared norm is
\begin{equation}\label{d-triangle-norm}
a+\frac r4.
\end{equation}
Indeed, if all squared coordinates are at most $1/2$, the half-triple
dominates them and is the average of the three tuples having one coordinate
equal to $1$ and the other two equal to $1/4$. Otherwise, a largest squared
coordinate $u\in[1/2,1]$ bounds the other two by $1/(4u)$. The resulting
linear functional is convex in $u$, so its maximum occurs at an endpoint.
Choosing the largest kernel coordinate for the unit coordinate proves
\eqref{d-triangle-norm}.

Fixing two copies of the point carrying $a$ in the fourfold maximum gives
six distinct sums and therefore
\begin{equation}\label{d-triangle-quartic}
T_F(k_w)
\geq a^4+a^3r+\frac23a^2r^2
\geq(a+r/4)^4.
\end{equation}
For the final inequality, expand the fourth power. Since $r\leq2a\leq3a$,
the quadratic and higher terms are at most
$(153/256)a^2r^2<(2/3)a^2r^2$ when $r>0$. The zero cases follow directly.
Supports of size at most two are covered by zero extension.

For a positive full kernel, a missing point $0,4,1$, or $2$ excludes the
sum $0,16,1$, or $14$, respectively, from the smaller fourfold sumset.
That sum has a positive contribution for the full kernel. This proves
strictness. Finally, the same pair-bound argument with all weights equal
to one gives total at most $3/2$ for any feasible three-point profile;
an additional progression constraint can only decrease this maximum.
\end{proof}

\subsection{Reduction of the feasible profiles}

An endpoint estimate will be used for maxima at the first and third
coordinates.

\begin{lemma}\label{d-endpoint-curves}
For $u\in[1/2,1]$, set
\begin{equation}\label{d-Jh}
J=\frac1{4u},\qquad h=J(1-J)^2.
\end{equation}
The curve $(u,J,h,J)$ is dominated coordinatewise by the convex hull of
$H_L,V_L,V_R$. The curve $(h,J,u,J)$ is dominated coordinatewise by the
convex hull of $H_R,V_R,V_L$.
\end{lemma}

\begin{proof}
Set
\begin{equation}\label{d-curve-weights}
\begin{gathered}
t=2u-1,\qquad s=\frac{(2u-1)(1-u)}{4u},\qquad \beta=\frac s2,\\
\alpha=t+\frac{23}{32}\beta,\qquad \gamma=1-\alpha-\beta.
\end{gathered}
\end{equation}
These weights are nonnegative and sum to one, since
\[
\gamma=(1-u)\left(2-\frac{55(2u-1)}{256u}\right)\geq0.
\]
In $\gamma H_L+\alpha V_L+\beta V_R$, the first coordinate equals $u$,
the second and fourth equal $J+201s/256$, and the third equals
\[
\frac{7+2u}{64}+\frac{1815s}{4096}.
\]
On the other hand,
\begin{equation}\label{d-curve-remainder}
h-\frac{7+2u}{64}
=s\frac{u^2+5u-1}{16u^2}\leq\frac{7s}{16},
\end{equation}
because $6u^2-5u+1=(3u-1)(2u-1)\geq0$. The inequality $1815>1792$
proves domination of the third coordinate. Interchanging the first and
third coordinates proves the second assertion.
\end{proof}

\begin{lemma}\label{d-reduction}
For every nonnegative $k$, its linear functional on the profiles satisfying
\eqref{d-feasibility} is bounded above by its largest value on the six
profiles in \eqref{d-profiles} and on the feasible profiles having a zero
coordinate.
\end{lemma}

\begin{proof}
First suppose all coordinates are at most $1/2$. Drop the second
progression constraint for an upper bound and increase $w_4$ to $1/2$.
For fixed $m=w_1\in[0,1/2]$, the first constraint gives
$w_0w_2\leq(1-m)^2/4$. Maximizing a nonnegative linear functional under
this product bound and the upper bounds $w_0,w_2\leq1/2$ reduces to
\[
(w_0,w_2)=\left(\frac12,\frac{(1-m)^2}{2}\right)
\quad\hbox{or its interchange}.
\]
To see this, on the product boundary the objective has the convex form
$a\ell+cK/\ell$; it is largest at an endpoint of the allowed interval.
The remaining objective is convex in $m$, so its endpoints give
$H_0,H_L,H_R$.

Otherwise choose a largest coordinate $u\geq1/2$. Each other coordinate
is at most $J=1/(4u)$. We consider its four possible locations.

\paragraph{Largest coordinate at index 1.}
Drop the second progression constraint and set $w_4=J$. The same
product-bound argument gives
\begin{equation}\label{d-middle-one}
(w_0,w_2)=(J,f)\quad\hbox{or its interchange},
\qquad f=u(1-u)^2.
\end{equation}
Assign the larger of $k(0),k(2)$ to $J$. The objective becomes
\[
k(1)u+k_{\rm small}(J+f)
+(k_{\rm large}-k_{\rm small}+k(4))J.
\]
It is convex in $u$, since $J''>0$ and
\[
(J+f)''=\frac1{2u^3}+6u-4>0.
\]
On $[1/2,2/3]$, the first summand is at least $27/16$ and the second
at least $-1$; above $2/3$ both are positive. The endpoint profiles are
$H_L,H_R$ and the feasible tuples
\[
(1/4,1,0,1/4),\qquad(0,1,1/4,1/4).
\]

\paragraph{Largest coordinate at index 0.}
Retain both progression constraints. For fixed $c=w_2\in[0,J]$, increase
the other two coordinates to
\begin{equation}\label{d-largest-zero-bounds}
w_1=\min\{J,1-\sqrt{c/J}\},\qquad w_4=J(1-c)^2.
\end{equation}
The first formula changes at $c=h$. The objective is convex on each of
$[0,h]$ and $[h,J]$: its nonaffine terms are a nonnegative quadratic
and, on the second interval, a nonnegative multiple of $-\sqrt c$.
The resulting endpoints are
\begin{equation}\label{d-largest-zero-profiles}
(u,J,0,J),\qquad
(u,J,h,J(1-h)^2),\qquad
(u,0,J,h).
\end{equation}
All three satisfy \eqref{d-feasibility}. The first and third have a zero
coordinate. The middle is dominated by $(u,J,h,J)$, to which
Lemma~\ref{d-endpoint-curves} applies.

\paragraph{Largest coordinate at index 2.}
Use $f=u(1-u)^2$ again. We have $f\leq h$: after multiplication by
$4u$, the equivalent inequality $1-J-2u(1-u)\geq0$ becomes
\[
(2u-1)(4u^2-2u+1)\geq0.
\]
For fixed $a=w_0>0$, the other admissible maxima are
\begin{equation}\label{d-largest-two-bounds}
w_1=\min\{J,1-2\sqrt{au}\},\qquad
w_4=\min\left\{J,\frac{(1-u)^2}{4a}\right\}.
\end{equation}
The objective is affine or convex on the intervals with endpoints
$0,f,h,J$, since the nonaffine terms are nonnegative multiples of
$1/a$ and $-\sqrt a$. On $[0,f]$ it increases with $a$. It is enough
to use the feasible profiles
\begin{equation}\label{d-largest-two-profiles}
(f,J,u,J),\qquad
\left(h,J,u,\frac{(1-u)^2}{4h}\right),\qquad
(J,0,u,f).
\end{equation}
When $u=1$, the condition $w_0w_4=0$ gives these same limiting endpoints
directly on its two branches. The last profile has a zero coordinate.
Since $h\geq f$, the middle one is dominated by $(h,J,u,J)$ and hence
by Lemma~\ref{d-endpoint-curves}. The first is dominated by the chord
from $H_R$ to the feasible tuple $(0,1/4,1,1/4)$, with weight $2u-1$
on the latter. Its first coordinate is $(1-u)/4\geq u(1-u)^2$;
its second and fourth dominate $J$ by convexity of $1/u$.

\paragraph{Largest coordinate at index 4.}
The second progression constraint gives
$w_0\leq J(1-w_2)^2$. Drop the first constraint for an upper bound.
The convex quadratic lies below its endpoint chord, so the profile is
dominated by a convex combination of
\begin{equation}\label{d-largest-four-curves}
B(u)=(J,J,0,u),\qquad W(u)=(J(1-J)^2,J,J,u).
\end{equation}
The weight on $W(u)$ is $w_2/J$. Both endpoints also satisfy the first
progression constraint; for $W(u)$ its left side is $3J-2J^2\leq1$
for $J\in[1/4,1/2]$.

Convexity of $1/u$ bounds $B(u)$ by the chord joining the feasible
tuples
\[
Q=(\tfrac12,\tfrac12,0,\tfrac12),\qquad
S=(\tfrac14,\tfrac14,0,1).
\]
For $W(u)$, put $Z=(1,0,1/4,9/64)$ and define
\begin{equation}\label{d-outer-weights}
\lambda=\frac{(1-u)(64u+55)}{87u},\qquad
\mu=\frac{(2u-1)(64u+23)}{87u},\qquad
\nu=\frac{32(2u-1)(1-u)}{87u}.
\end{equation}
The weights are nonnegative and sum to one. The tuple
$\lambda H_R+\mu R+\nu Z$ has coordinates $w_1=J$, $w_4=u$,
and $w_2=J+\nu/4$. Its first coordinate exceeds $J(1-J)^2$ by
\begin{equation}\label{d-outer-remainder}
\frac{(1-u)(2u-1)(576u^2-145u+29)}{1856u^3}\geq0.
\end{equation}
The quadratic is positive because its leading coefficient is positive
and its discriminant is $21025-66816<0$. The profiles $Z,Q,S$ are
feasible and have a zero coordinate. This finishes all four cases,
including ties.
\end{proof}

\subsection{Quartic comparisons}

Write $k=(A,B,C,D)$ on $(0,1,2,4)$. In the expansion of $(k\cdot v)^4$,
group monomials according to the sum of their four support indices.
Each group is at most its coefficient sum times the maximum defining
$T_F(k)$ at that index. Thus coefficient sums at most one prove the
desired comparison. The following computations use weighted AM--GM
to move a term to other fourfold products when a coefficient sum
exceeds one.

\begin{lemma}\label{d-half-comparisons}
For $v=H_0,H_L,H_R,R$ and $k\geq0$,
\begin{equation}\label{d-half-inequality}
(k\cdot v)^4\leq T_F(k).
\end{equation}
Each inequality is strict when $k$ is positive on $F$.
\end{lemma}

\begin{proof}
For $H_0$, the target is $(A+C+D)^4/16$ on the progression
$\{0,2,4\}$. After dividing indices by two, its coefficient sums are
\[
\frac1{16}(1,4,10,16,19,16,10,4,1).
\]
Use $A^2D^2\leq(A^3D+AD^3)/2$ on the entire term
$(6/16)A^2D^2$. The coefficient sums at indices $2,4,6$ become
$13/16$ each; all others were at most one. This compares with the
three-point maximum. The positive full kernel also has a contribution
at the missing odd sum $1$, giving strictness.

For $H_L$, its index polynomial is
$(1+z+z^2/4+z^4)/2$. If $c_n$ is the coefficient of its fourth power,
the complete coefficient sums are
\begin{equation}\label{d-HL-coefficients}
\begin{gathered}
\begin{array}{c rrrrrrrrr}
\toprule
n&0&1&2&3&4&5&6&7&8\\
\midrule
4096c_n&256&1024&1792&1792&2144&3520&3952&2576&2497\\
\bottomrule
\end{array}\\[6pt]
\begin{array}{c rrrrrrrr}
\toprule
n&9&10&11&12&13&14&15&16\\
\midrule
4096c_n&3264&2320&768&1120&1024&256&0&256\\
\bottomrule
\end{array}
\end{gathered}
\end{equation}
These follow by squaring the row
$(1,2,3/2,1/2,33/16,2,1/2,0,1)$ and dividing by $16$.
Every coefficient is at most $247/256<1$.

For $H_R$, the linear form is $A/8+(B+C+D)/2$. Its coefficient
sums $c_n$ are
\begin{equation}\label{d-HR-coefficients}
\begin{gathered}
\begin{array}{c rrrrrrrrr}
\toprule
n&0&1&2&3&4&5&6&7&8\\
\midrule
256c_n&1/16&1&7&28&71&124&172&224&262\\
\bottomrule
\end{array}\\[6pt]
\begin{array}{c rrrrrrrr}
\toprule
n&9&10&11&12&13&14&15&16\\
\midrule
256c_n&240&208&192&112&64&64&0&16\\
\bottomrule
\end{array}
\end{gathered}
\end{equation}
Only sum $8$ exceeds one. Replace its term $(3/128)A^2D^2$ by
$(3/256)A^3D+(3/256)AD^3$. The coefficient sums at indices $4,8,12$
become $74/256,1,115/256$, respectively. All others pass unchanged.
The coefficient $1/4096$ at index $0$ gives strictness when $A>0$.

For $R$, relabel the support by $p\mapsto4-p$. Its index polynomial becomes
\[
C(z)=1+\tfrac14z^2+\tfrac14z^3+\tfrac9{64}z^4.
\]
The coefficients of $C(z)^4$ at indices $0$ through $7$ are
\[
1,\quad0,\quad1,\quad1,\quad\tfrac{15}{16},\quad
\tfrac34,\quad\tfrac{55}{64},\quad\tfrac{39}{64}.
\]
For indices at least $8$, write $C=1+H$. In
$C^4=1+4H+6H^2+4H^3+H^4$, the first two terms do not contribute.
The contribution of $6H^2$ is at most $3/8$ and occurs only at index $8$.
Since $H$ is coefficientwise bounded by $z^2(1+z+z^2)/4$, the other
two contributions are at most $7/16$ and $19/256$. Here $7$ and $19$
are the largest coefficients of $(1+z+z^2)^3$ and $(1+z+z^2)^4$.
The total is at most $227/256<1$. Thus all coefficients pass, and the
coefficient $15/16$ at reflected index $4$ gives strictness for positive
kernels. Reflection reverses the output indices and preserves their sum.
\end{proof}

\begin{lemma}\label{d-VL-comparison}
For $k\geq0$, $(k\cdot V_L)^4\leq T_F(k)$, strictly when $k$ is positive
on $F$.
\end{lemma}

\begin{proof}
The linear form is $E=A+B/4+9C/64+D/4$. The coefficient sums in $E^4$
are listed below. An omitted entry is zero.
\begin{equation}\label{d-VL-ledger}
\begin{array}{r l r l}
\toprule
n&c_n&n&c_n\\
\midrule
0&1&9&3/16+243/16384\\
1&1&10&33/256+729/262144\\
2&15/16&11&27/1024\\
3&31/64&12&1/16+243/32768\\
4&1+467/2048&13&1/64\\
5&3/4+279/4096&14&9/1024\\
6&39/64+1215/65536&15&0\\
7&29/128+729/262144&16&1/256\\
8&3/8+351/4096+6561/16777216&&\\
\bottomrule
\end{array}
\end{equation}
Only sum $4$ exceeds one. Its polynomial is
\[
A^3D+\frac{243}{2048}A^2C^2+\frac{27}{256}AB^2C+\frac1{256}B^4
\leq U+aV+bW,
\]
where
\begin{equation}\label{d-VL-collision}
\begin{gathered}
U=A^3D,\quad V=A^2C^2,\quad W=B^4,\\
a=\frac{351}{2048},\qquad b=\frac{29}{512},\qquad
a+b=\frac{467}{2048}.
\end{gathered}
\end{equation}
This uses $AB^2C\leq(A^2C^2+B^4)/2$. For $M=\max(U,V,W)$,
\begin{equation}\label{d-minimum-split}
U+aV+bW\leq M+a\min(U,V)+b\min(U,W).
\end{equation}
Indeed,
$a(V-U)_++b(W-U)_+\leq(a+b)(M-U)\leq M-U$.
Weighted AM--GM also gives
\begin{equation}\label{d-scaled-minima}
\begin{aligned}
\min(A^3D,A^2C^2)&\leq\tfrac4{15}A^3C+\tfrac{25}{12}A^2D^2,\\
\min(A^3D,B^4)&\leq\tfrac4{15}A^2B^2+\tfrac{25}{12}A^2D^2.
\end{aligned}
\end{equation}
For the first line, the minimum is at most $U^{2/3}V^{1/3}$, which equals
$(A^3C)^{2/3}(A^2D^2)^{1/3}$. The stated bound follows since
$(4/15)^2(25/12)=4/27$. The second line has the same exponents.

The term $M$ is at most the maximum at sum $4$. All added terms from
\eqref{d-scaled-minima} go to sums $2$ and $8$. The final coefficient
sum at $2$ is
\begin{equation}\label{d-VL-two}
\frac{15}{16}+\frac{467}{2048}\frac4{15}
=\frac{7667}{7680}<1.
\end{equation}
The original coefficient sum at $8$ is less than $15/32$, and the added
weight is
\[
\frac{467}{2048}\frac{25}{12}=\frac{11675}{24576}<\frac12.
\]
Its final total is below $31/32$. The remaining sums pass by
\eqref{d-VL-ledger}. For positive kernels, the positive maximum and
remaining slack at sum $2$ give strictness.
\end{proof}

\begin{lemma}\label{d-VR-comparison}
For $k\geq0$, $(k\cdot V_R)^4\leq T_F(k)$, strictly when $k$ is positive
on $F$.
\end{lemma}

\begin{proof}
Relabel the support by $p\mapsto2-p$, and label the kernel values on positions
$(0,1,2,-2)$ by $(A,B,C,D)$. The profile is then
$(1,1/4,9/64,1/4)$, so the target is again $E=A+B/4+9C/64+D/4$.
Only sums $0,1,2,-2$ have coefficient sum greater than one.
Retain respectively
\begin{equation}\label{d-VR-retained}
A^4,\qquad A^3B,\qquad
\tfrac9{16}A^3C+\tfrac38A^2B^2,\qquad A^3D.
\end{equation}
The terms removed from these groups are
\begin{equation}\label{d-VR-removed}
\begin{array}{r l}
\toprule
\text{sum}&\text{removed terms}\\
\midrule
0&(3/16)AB^2D+(27/64)A^2CD+(243/32768)C^2D^2\\
1&(1/64)B^3D+(27/128)ABCD\\
2&(27/1024)B^2CD+(243/4096)AC^2D\\
-2&(3/128)B^2D^2+(27/256)ACD^2\\
\bottomrule
\end{array}
\end{equation}
Apply the following weighted AM--GM inequalities:
\begin{equation}\label{d-VR-redistribution}
\begin{aligned}
AB^2D&\leq(A^2D^2+B^4)/2,&
A^2CD&\leq(A^2C^2+A^2D^2)/2,\\
C^2D^2&\leq(C^4+D^4)/2,&
B^3D&\leq(3B^4+D^4)/4,\\
ABCD&\leq(A^2D^2+B^2C^2)/2,&
B^2CD&\leq B^4/2+C^4/4+D^4/4,\\
AC^2D&\leq(A^2D^2+C^4)/2,&
B^2D^2&\leq(B^4+D^4)/2,\\
ACD^2&\leq(A^2C^2+D^4)/2.&&
\end{aligned}
\end{equation}
Every destination is a fourfold product on the reflected support. Only
sums $4,-4,6,8,-8$ receive additions. Their coefficient sums are
\begin{equation}\label{d-VR-ledger}
\begin{array}{r l l l}
\toprule
\text{sum}&\text{original}&\text{addition}&\text{final}\\
\midrule
4&467/2048+729/262144&807/2048&163801/262144<5/8\\
-4&3/8+9/1024&3603/8192&6747/8192<7/8\\
6&1215/65536&27/256&8127/65536<1/8\\
8&6561/16777216&2619/65536&677025/16777216<1/16\\
-8&1/256&5155/65536&5411/65536<3/32\\
\bottomrule
\end{array}
\end{equation}
Besides these and the four groups in \eqref{d-VR-retained}, the
unchanged nonzero coefficient sums are
\begin{equation}\label{d-VR-unchanged}
\begin{array}{r l r l}
\toprule
\text{sum}&\text{coefficient sum}&\text{sum}&\text{coefficient sum}\\
\midrule
-6&1/16&3&31/64+243/16384\\
-5&1/64&5&279/4096\\
-3&3/16&7&729/262144\\
-1&3/4+27/1024&&\\
\bottomrule
\end{array}
\end{equation}
All are less than one. The retained groups at $0,1,-2$ have coefficient
sum one, and the group at $2$ has coefficient sum $15/16$. This proves
the inequality and gives strictness for positive kernels at sum $2$.
Reflection preserves the sum of the fourfold maxima.
\end{proof}

\begin{proposition}\label{d-certificate}
For $F=\{0,1,2,4\}$ and every nonnegative kernel $k$,
\begin{equation}\label{d-main-inequality}
N_F(k)^4\leq T_F(k).
\end{equation}
The inequality is strict if $k$ is positive on $F$.
\end{proposition}

\begin{proof}
Lemma~\ref{d-reduction} bounds the weighted objective by its values
on \eqref{d-profiles} and on feasible profiles with a zero coordinate.
Lemmas~\ref{d-boundary}, \ref{d-half-comparisons},
\ref{d-VL-comparison}, and~\ref{d-VR-comparison} bound the fourth power
of each of these values by $T_F(k)$.

For positive kernels every listed comparison is strict. The feasible
set is compact by \eqref{d-feasibility}, so its weighted maximum is
attained. Applying the strict comparison to a maximizing profile
proves the strict form of \eqref{d-main-inequality}.
\end{proof}

\subsection{The sharp unweighted norm}

The same reduction determines the unweighted norm exactly.

\begin{proposition}\label{d-sharp-norm}
For $F=\{0,1,2,4\}$,
\begin{equation}\label{d-norm-value}
\sup_{\substack{0\ne x\geq0\\ \operatorname{supp}(x)\subseteq F}}
\frac{\|x\|_2^2}{\|x*x\|_\infty}=\frac{105}{64}.
\end{equation}
Equality holds for $x=(3/8,1/2,1/2,1)$ in the coordinate order
$(0,1,2,4)$.
\end{proposition}

\begin{proof}
By homogeneity, normalize $\|x*x\|_\infty=1$. In
Lemma~\ref{d-reduction}, the totals of $V_L,V_R,R$ are $105/64$,
those of $H_L,H_R$ are $13/8=104/64$, and that of $H_0$ is $3/2$.
Every feasible profile with a zero coordinate has total at most $3/2$
by Lemma~\ref{d-boundary}. This gives the upper bound.

The squared profile of the stated $x$ is $R$. It satisfies
\eqref{d-feasibility}: the two progression left sides are $5/8$ and
$1$, all pair bounds hold, and the coordinate at $4$ is $1$.
Thus $\|x*x\|_\infty=1$ and $\|x\|_2^2=105/64$.
\end{proof}

\section{The support-size boundary}
\label{sec:boundary}

\begin{proposition}\label{b-small-support}
For every nonempty integer set $F$ with $|F|\le3$ and every nonnegative
kernel $k$ supported on $F$, one has $N_F(k)^4\le T_F(k)$.
\end{proposition}

\begin{proof}
Sets with at most two points are Sidon, as is any three-point set that is
not an arithmetic progression.  Proposition~\ref{c-sidon} applies to all
of them.  For a three-term progression, complete it to a four-term
progression and extend both $x$ and $k$ by zero.  This preserves ordinary
convolution, the weighted objective, and $T_F(k)$.  The four-term-progression
inequality proved above therefore gives the desired bound for every feasible
$x$, and then for $N_F(k)$.
\end{proof}

For a singleton $F=\{p\}$ the bound is equality:
$N_F(k)=k(p)$ and $T_F(k)=k(p)^4$.  The corresponding max-convolution
inequality is also equality for every nonzero $f,g$ allowed in
Corollary~\ref{cor:h4}, since
\[
 \|f\star g\|_1=g(p)\|f\|_1,\qquad
 \|f\star\widetilde{gk}\|_1=g(p)k(p)\|f\|_1.
\]
Thus the strict four-point statement does not extend to singleton supports.

The preceding proposition, the classification in
Proposition~\ref{c-classification}, and the five four-point cases prove
the positive part of Theorem~\ref{thm:main}.  The following example proves
its support-size assertion.

\begin{proposition}\label{b-five-point}
Let $F=\{0,2,7,8,11\}$ and $k=1_F$.  Then
\[
 N_F(k)=\frac52,\qquad T_F(k)=39,\qquad
 N_F(k)^4-T_F(k)=\frac1{16}.
\]
\end{proposition}

\begin{proof}
The fifteen unordered pair sums, including diagonal pairs, are
\[
 0,2,4,7,8,9,10,11,13,14,15,16,18,19,22.
\]
They are distinct, so $F$ is Sidon.  For $x=1_F/\sqrt2$, ordinary
convolution equals $1/2$ at diagonal sums and $1$ at off-diagonal sums.
Consequently $x*x\le1$ and $\sum_{p\in F}x(p)^2=5/2$.
Equation~\eqref{c-sidon-norm}, which holds for Sidon sets of any size,
also supplies the upper bound:
$N_F(1_F)=\max\{5/2,(5+3)/4\}=5/2$.

To count $4F$, put $E=\{0,2,7,8\}$ and write
$[r,s]_{\mathbb Z}=\{r,r+1,\ldots,s\}$ for integer endpoints $r\le s$.
For $h\ge0$, partition the summands in $hE$ by the number $m$ drawn
from $\{7,8\}$.  Their sum ranges from $7m$ to $8m$ in unit steps,
while the remaining summands from $\{0,2\}$ contribute $2j$ for
$0\le j\le h-m$.  Hence
\begin{equation}\label{b-small-sumsets}
 hE=\bigcup_{m=0}^{h}\ \bigcup_{j=0}^{h-m}
           [7m+2j,8m+2j]_{\mathbb Z}.
\end{equation}
In particular,
\[
 4E=\{0,2,4\}\cup[6,26]_{\mathbb Z}\cup[28,32]_{\mathbb Z}.
\]
Now partition $4F$ by the number $\ell$ of summands equal to $11$.
Formula~\eqref{b-small-sumsets} gives the new elements at each step:
\begin{center}
\begin{tabular}{clc}
\toprule
Layer & Elements not in preceding layers & Number \\
\midrule
$4E$ & $\{0,2,4\}\cup[6,26]_{\mathbb Z}\cup[28,32]_{\mathbb Z}$ & $29$ \\
$11+3E$ & $\{27,33,34,35\}$ & $4$ \\
$22+2E$ & $\{36,37,38\}$ & $3$ \\
$33+E$ & $\{40,41\}$ & $2$ \\
$\{44\}$ & $\{44\}$ & $1$ \\
\bottomrule
\end{tabular}
\end{center}
Thus
\[
 4F=\{0,2,4\}\cup[6,38]_{\mathbb Z}\cup\{40,41,44\},
 \qquad |4F|=39.
\]
The fourfold max-convolution of an indicator is the indicator of its
fourfold sumset, so $T_F(1_F)=39$.  Finally,
$(5/2)^4=625/16=39+1/16$.
\end{proof}

The example shows failure of the sufficient covering bound
$N_F(k)^4\le T_F(k)$ with both $x$ and $k$ positive throughout a
five-point set.  A counterexample to the reflected max-convolution
inequality would additionally require functions $f,g$ for which the ratio
in \eqref{eq:ratio} exceeds $T_F(k)^{1/4}$; this construction supplies
no such functions.

\section{Sharp unweighted covering constants}
\label{sec:constants}

The five pair-sum patterns have different unweighted norms.  Their values,
with one maximizing squared profile in the natural coordinate order, are
listed below.
\begin{center}
\begin{tabular}{lcl}
\toprule
Four-point pattern & $N_F(1_F)$ & Maximizing squared profile \\
\midrule
Sidon & $2$ & $(1/2,1/2,1/2,1/2)$ \\
Isolated progression, $(0,a,2a,b)$ & $7/4$ & $(1/4,1/4,1/4,1)$ \\
Two progressions, $(0,1,2,4)$ & $105/64$ & $(9/64,1/4,1/4,1)$ \\
Parallelogram, $(0,a,b,a+b)$ & $25/16$ & $(1,1/4,1/4,1/16)$ \\
Four-term progression & $3/2$ & $(1,1/4,0,1/4)$ \\
\bottomrule
\end{tabular}
\end{center}

For Sidon supports, \eqref{c-sidon-norm} gives
$\max\{2,7/4\}=2$.  The isolated-progression profile sums are
$105/64$ twice, $3/2$ four times, and $7/4$ twice.  The sixteen
parallelogram profile sums are $3/2$ for the twelve profiles supported
on three corners and $25/16$ for the four tensor profiles.  For a
four-term progression, the five profile sums before reversal are
$3/2,3/2,381/256,3/2,89/64$.  The double-progression reduction bounds
every profile by ones with total at most $105/64$, and its displayed
profile is feasible.  Thus the reductions prove all five upper bounds,
while the listed feasible profiles attain them.

By Lemma~\ref{lem:cover}, each value is also the maximum of
$L_F(g,1_F)$ over nonzero nonnegative $g$, and the supremum when $g$
is required to be positive at every point.  This sharpness concerns the
universal translated covering problem.  The max-convolution ratio in
\eqref{eq:ratio} can be smaller than its covering bound.

\section{Constant kernels on five points}
\label{sec:five-point}

The reflected inequality extends beyond the support-size boundary of the
covering bound when the kernel is constant.

\begin{theorem}\label{thm:five-point}
Let $F\subset\mathbb Z$ have five elements. For every nonzero finitely
supported $f\ge0$ and every nonzero $g\ge0$ supported on $F$,
\[
 \|f\star\widetilde g\|_1^4
 < |4F|\,\|f\star g\|_1^4.
\]
In particular, the five-point set in Proposition~\ref{b-five-point}
satisfies this inequality despite failing the sufficient covering bound.
\end{theorem}

We first prove a bound that applies to every five-point support, then
separate Sidon supports from those with pair-sum collisions.

\subsection{A bound from energy and the range of the weights}

\begin{lemma}\label{five-ratio}
For every five-element $F\subset\mathbb Z$ and every $f,g$ as in
Theorem~\ref{thm:five-point},
\[
 \frac{\|f\star\widetilde g\|_1}{\|f\star g\|_1}\le\frac{62}{25}.
\]
Consequently, Theorem~\ref{thm:five-point} holds whenever $|4F|\ge38$.
\end{lemma}
\begin{proof}
For a nonempty finite set $A$, let $r_+$ and $r_-$ count the representations
of a point as $a+p$ and $a-p$, respectively, with $(a,p)\in A\times F$.
Put $P=|A+F|$, $Q=|A-F|$ and $W=5|A|$. The two quadratic energies are
equal: exchanging the two $F$ coordinates bijects the solutions of
$a+p=a'+p'$ with those of $a-p=a'-p'$. Write their common value as $E$.
Since $1\le r_-(s)\le5$ on its support, $r_-(s)^2\le6r_-(s)-5$.
Cauchy--Schwarz applied to $r_+$ gives
\[
 \frac{W^2}{P}\le E\le6W-5Q,
 \qquad
 \frac QP\le\frac{6t-t^2}{5}\le\frac95,
 \quad t=\frac WP.
\]
Integrate this inequality over the level sets of $f$ to obtain
\begin{equation}\label{five-energy}
 \|f\star1_{-F}\|_1\le\frac95\|f\star1_F\|_1.
\end{equation}
For $g>0$ on $F$, put $r=\max_Fg/\min_Fg\ge1$. Pointwise comparison
with the indicator kernels gives
\begin{equation}\label{five-range-energy}
 \frac{\|f\star\widetilde g\|_1}{\|f\star g\|_1}\le\frac95r.
\end{equation}

The dual in Lemma~\ref{lem:cover}, with $k=1_F$, has objective
$S=\sum_iw_i$, where $w_i=g_i z_i$ and $g_i=g(p_i)>0$ for
$F=\{p_1,\ldots,p_5\}$. Its constraints imply
\begin{equation}\label{five-dual-pairs}
 0\le w_i\le1,\qquad
 \frac{g_j}{g_i}w_i+\frac{g_i}{g_j}w_j\le1,\qquad
 w_iw_j\le\frac14\quad(i\ne j).
\end{equation}
The middle inequality retains two nonnegative terms of the constraint
at $p_i+p_j$; AM--GM gives the last one.

Every four coordinates have sum at most two. If their largest value
$u$ is at most $1/2$, this is immediate. Otherwise their sum is at most
$u+3/(4u)\le2$ for $1/2\le u\le1$. It follows that $w_i\ge S-2$
for every $i$. Choose distinct indices at which $g$ takes its maximum
and minimum; when $r=1$, any two distinct indices suffice. Their
constraint in \eqref{five-dual-pairs} gives
\[
 1\ge(r+r^{-1})(S-2),\qquad
 S\le2+\frac1{r+r^{-1}}.
\]
Together with the covering lemma, this proves
\begin{equation}\label{five-range-cover}
 \frac{\|f\star\widetilde g\|_1}{\|f\star g\|_1}
 \le\min\left\{\frac95r,\ 2+\frac1{r+r^{-1}}\right\}.
\end{equation}
For $1\le r\le4/3$, the first bound is at most $12/5$.
For $r\ge4/3$, the second is at most $2+12/25=62/25$.

For nonnegative nonzero $g$, apply the bound to $g+\epsilon1_F$ and
let $\epsilon\downarrow0$. For fixed $f$, both max-convolutions are
supported in fixed finite sets and their values are continuous in $\epsilon$.
The limiting denominator is positive. Finally,
\[
 \left(\frac{62}{25}\right)^4
 =\frac{14776336}{390625}<38,
\]
which proves the last assertion.
\end{proof}

\subsection{Supports with a pair-sum collision}

For an ordinary-autoconvolution-feasible $x\ge0$, write $w_p=x(p)^2$.
The universal constraints are $0\le w_p\le1$ and
$w_pw_q\le1/4$ for distinct points. We use repeatedly the elementary
inequality
\begin{equation}\label{five-box}
 b+c\le a+\frac{bc}{a}\qquad(0\le b,c\le a, a>0),
\end{equation}
which follows from $(a-b)(a-c)\ge0$.

\begin{lemma}\label{five-parallelogram}
If a five-element integer set $F$ contains four distinct points forming a
parallelogram, then $N_F(1_F)\le2$.
\end{lemma}
\begin{proof}
Label the four points so that $p_0+p_3=p_1+p_2$, and label the fifth
point $p_4$. The extra convolution constraint is
\[
 \sqrt{w_0w_3}+\sqrt{w_1w_2}\le\frac12.
\]
If all weights are at most $1/2$, apply \eqref{five-box} to the two
opposite pairs. With $p=\sqrt{w_0w_3}$ and $q=\sqrt{w_1w_2}$,
\[
 w_0+w_1+w_2+w_3\le1+2(p^2+q^2)
 \le1+2(p+q)^2\le\frac32.
\]
Adding $w_4\le1/2$ proves the bound in this region.

Otherwise let $u\in[1/2,1]$ be a largest weight, and put
$a=1/(4u)\in[1/4,1/2]$. Every other weight is at most $a$.
If $u=w_4$, then
\[
 \sum_iw_i\le u+2a+\frac{p^2+q^2}{a},
 \qquad 0\le p,q\le a,\quad p+q\le\frac12.
\]
On this polygon the maximum of $p^2+q^2$ is
$a^2+(1/2-a)^2$, as follows by checking its vertices. Hence
\[
 \sum_iw_i\le2u+\frac1u-1\le2.
\]
The last function is convex and equals two at both endpoints.

If the largest weight is a corner, relabel so that $u=w_0$ and put
$p=\sqrt{w_1w_2}\in[0,a]$. The opposite weight is at most
$(1/2-p)^2/u$. Equation~\eqref{five-box} and $w_4\le a$ give
\[
 \sum_iw_i\le u+2a+\frac{p^2}{a}+\frac{(1/2-p)^2}{u}.
\]
This is convex in $p$, so the endpoint bounds are
\[
 A(u)=u+\frac{3}{4u},\qquad
 B(u)=u+\frac1u-\frac1{4u^2}+\frac1{16u^3}.
\]
Both are convex on $[1/2,1]$; for the second,
$B''(u)=(2u^2-3u/2+3/4)/u^5>0$.
Their endpoint values are $(2,7/4)$ and $(2,29/16)$, respectively.
Thus every feasible profile has total at most two.
\end{proof}

\begin{lemma}\label{five-progression}
If a five-element integer set $F$ contains a three-term progression,
then $N_F(1_F)\le17/8$.
\end{lemma}
\begin{proof}
Write $b,c$ for the endpoint weights and $v$ for the middle weight.
The extra constraint is
\begin{equation}\label{five-ap-constraint}
 v+2\sqrt{bc}\le1.
\end{equation}
If every weight is at most $1/2$, then \eqref{five-box} gives
\[
 b+c+v\le\frac12+2bc+v
 \le\frac12+\frac{(1-v)^2}{2}+v
 =1+\frac{v^2}{2}\le\frac98.
\]
The other two weights contribute at most one.

Otherwise let $u\in[1/2,1]$ be a largest weight and set $a=1/(4u)$.
All remaining weights are at most $a$. There are three cases.

\paragraph{An endpoint is largest.}
Take $b=u$. For $v\in[0,a]$, the other endpoint is at most
$(1-v)^2/(4u)\le a$, so
\[
 \sum w\le u+2a+v+\frac{(1-v)^2}{4u}.
\]
Convexity in $v$ reduces this to $v=0,a$, giving $A(u)$ above and
\[
 C(u)=u+\frac1u-\frac1{8u^2}+\frac1{64u^3}.
\]
The function $C$ is convex because
$C''(u)=(2u^2-3u/4+3/16)/u^5>0$.
Its endpoint values are $17/8$ and $121/64$; also $A(u)\le2$.

\paragraph{The middle point is largest.}
Take $v=u$. Equations~\eqref{five-box} and \eqref{five-ap-constraint}
give $b+c\le a+u(1-u)^2$. Thus
\[
 \sum w\le D(u)=u+\frac{3}{4u}+u(1-u)^2.
\]
Here $D''(u)=6u-4+3/(2u^3)>0$ on $[1/2,1]$, and its endpoint
values are $17/8$ and $7/4$.

\paragraph{An outside point is largest.}
The other outside weight is at most $a$, and
\[
 b+c+v\le v+a+\min\left\{a,\frac{(1-v)^2}{4a}\right\},
 \qquad 0\le v\le a.
\]
If $u\in[3/4,1]$, the simple bound $b+c+v\le3a$ gives
$\sum w\le u+1/u\le25/12<17/8$.
If $u\in[1/2,3/4]$, then $a\in[1/3,1/2]$. The two branches
meet at $v_0=1-2a\in[0,a]$. The first branch increases to one.
The second is convex, with endpoint values one and
$2a+(1-a)^2/(4a)$. Therefore
\[
 \sum w\le\max\left\{u+\frac1{4u}+1,\
                       2u+\frac{13}{16u}-\frac12\right\}.
\]
Both functions are convex on $[1/2,3/4]$. Their endpoint values are
$(2,25/12)$ and $(17/8,25/12)$, respectively. This completes the cases.
\end{proof}

\begin{proposition}\label{five-nonsidon}
If a five-element integer set $F$ is not Sidon, then
$N_F(1_F)^4<|4F|$.
\end{proposition}
\begin{proof}
If $F$ contains a four-point parallelogram, Lemma~\ref{five-parallelogram}
gives $N_F(1_F)^4\le16<17\le|4F|$. The last bound follows from
repeated application of $|A+B|\ge|A|+|B|-1$ in $\mathbb Z$.

Otherwise every pair-sum collision consists of one diagonal pair and one
off-diagonal pair, forming a three-term progression. A sum cannot have two
different off-diagonal pairs, since they would be disjoint and form a
parallelogram. Each point is the centre of at most one progression, and
the two extreme points are never centres. Thus there are at most three
collisions among the fifteen unordered pair sums, giving
\[
 |2F|\ge12,\qquad |4F|\ge2|2F|-1\ge23.
\]
Lemma~\ref{five-progression} now gives
$N_F(1_F)^4\le83521/4096<23$.
\end{proof}

\subsection{Five-point Sidon supports}

\begin{proposition}\label{five-sidon-size}
Every five-element Sidon subset $F$ of $\mathbb Z$ satisfies $|4F|\ge39$.
Equality is attained by $F=\{0,2,7,8,11\}$.
\end{proposition}
\begin{proof}
We use Freiman's classical $3k-4$ theorem in the following form:
if $A\subset\mathbb Z$, $|A|=k\ge3$ and $|A+A|\le3k-4$, then $A$
is contained in an arithmetic progression with at most
$|A+A|-k+1$ elements. This is the equal-set specialization of the
statement recalled by Bollob\'as, Leader and Tiba~\cite[p.~1]{blt}.

Translate $F$ and divide by the greatest common divisor of its
differences, writing
\[
 F=\{0,a,b,c,d\},\qquad 0<a<b<c<d,\qquad \gcd(F)=1.
\]
These operations preserve the Sidon property and sumset cardinalities.
Set $B=F+F$. Then $|B|=15$, $\min B=0$, $\max B=2d$, and
$\gcd(B)=1$ because $F\subset B$.
If $|4F|=|B+B|\le38$, the theorem applies, since $38\le3\cdot15-4$.
It puts $B$ in an arithmetic progression with at most $38-15+1=24$
elements. Its step divides every difference of $B$, so its absolute
value is one. Consequently $2d\le23$, and $d\le11$.

The ten positive differences of a five-point Sidon set are distinct,
so $d\ge10$. If $d=10$, those differences would be $1,\ldots,10$,
whose sum is odd. Their sum is instead $4d+2c-2a$, which is even.
Thus the supposition $|4F|\le38$ forces $d=11$.

We classify this last case by hand. The ten differences omit one number
$m$ from $1,\ldots,11$, and the parity argument shows that $m$ is even.
Put $P_i=\{i,11-i\}$ for $1\le i\le5$. The six endpoint differences
\[
 a,b,c,11-a,11-b,11-c
\]
occupy exactly three distinct pairs $P_i$. The three remaining
differences $b-a,c-b,c-a$ form a Schur triple: the sum of the smaller
two is the largest. They must be three of the four members of the two
unused pairs, say $P_r,P_s$ with $r<s$.

Deleting one element of the ordered list $r,s,11-s,11-r$ must therefore
leave a Schur triple. Deleting the smallest cannot work, since
$s+(11-s)=11>11-r$. The other deletions give $s=2r$, $2r+s=11$, or
$r+2s=11$. Their complete solution list is
\begin{center}
\begin{tabular}{cccc}
\toprule
Omitted $m$ & Unused $(r,s)$ & Interior differences & $c-a$\\
\midrule
$2$  & $(1,2)$ & $\{1,9,10\}$ & $10$\\
$4$  & $(2,4)$ & $\{2,7,9\}$  & $9$\\
$6$  & $(3,5)$ & $\{3,5,8\}$  & $8$\\
$8$  & $(3,4)$ & $\{3,4,7\}$  & $7$\\
$10$ & $(1,5)$ & $\{1,5,6\}$  & $6$\\
\bottomrule
\end{tabular}
\end{center}
The interior points must occupy the other three complementary pairs,
once each. For $m=2$, the span $c-a=10$ is impossible inside $[1,10]$.
For $m=4$, it forces $(a,c)=(1,10)$, using the same complementary pair.
For $m=6$, the possibilities $(a,c)=(1,9),(2,10)$ give only
$b=4,7$, respectively. For $m=8$, the possible endpoint pairs are
$(1,8),(2,9),(3,10)$; the outer two use the forbidden $P_3$, and the
middle one repeats $P_2$. For $m=10$, the choices $a=1,4$ use the
forbidden $P_1$. The remaining choices $(a,c)=(2,8),(3,9)$ give only
$b=7,4$, respectively. Up to reflection, the only sets are therefore
\[
 F_1=\{0,1,4,9,11\},\qquad F_2=\{0,2,7,8,11\}.
\]

For the first representative,
\[
 2F_1=\{0,1,2,4,5,8,9,10,11,12,13,15,18,20,22\}.
\]
Successive elements are at most three apart, so
$2F_1+\{0,1,2\}$ contains $[0,24]$. The rest of $[0,38]$ is supplied by
\[
\begin{aligned}
\relax[25,28]&=15+[10,13],& [29,33]&=20+[9,13],\\
 [34,35]&=22+[12,13],& 36&=18+18,\\
 37&=15+22,&38&=18+20.
\end{aligned}
\]
All summands belong to $2F_1$. Also $40=18+22$, $42=20+22$ and
$44=22+22$. The only other candidates are $39,41,43$; each requires a
summand $20$ or $22$, but its complementary summand is absent. Hence
\[
 4F_1=[0,38]\cup\{40,42,44\},\qquad |4F_1|=42.
\]
For the second representative, Proposition~\ref{b-five-point} gives
\[
 4F_2=\{0,2,4\}\cup[6,38]\cup\{40,41,44\},\qquad |4F_2|=39.
\]
Here all intervals are integer intervals. Both possibilities contradict
$|4F|\le38$, and the second proves sharpness.
\end{proof}

The proof classifies normalized span $11$ only. It does not classify all
equality cases for $|4F|=39$.


\begin{proof}[Proof of Theorem~\ref{thm:five-point}]
For non-Sidon supports, combine Proposition~\ref{five-nonsidon} with
Lemma~\ref{lem:cover}. For Sidon supports, use
Proposition~\ref{five-sidon-size} and Lemma~\ref{five-ratio}.
\end{proof}

For $f=1_A$ and $g=1_F$, Theorem~\ref{thm:five-point} gives
$|A-F|^4<|A+F|^4|4F|$ for every finite nonempty integer set $A$ and every
five-element integer set $F$. The theorem also permits arbitrary
nonnegative weights in $f$ and $g$; the kernel $k$ is fixed to $1_F$.

\section{Arbitrary kernels on a five-point Sidon support}
\label{sec:five-point-weighted}

The support that defeats the universal covering bound also satisfies the
reflected inequality for every choice of kernel weights. The proof uses
rational coefficient certificates in addition to the estimates below.

\begin{theorem}\label{thm:five-point-weighted}
Let $F=\{0,2,7,8,11\}$. Every nonnegative kernel $k$ supported on $F$
and every pair of nonnegative finitely supported functions $f,g$ on
$\mathbb Z$ satisfy
\[
 \|f\star\widetilde{gk}\|_1^4
 \le T_F(k)\,\|f\star g\|_1^4.
\]
The inequality is strict when $k$ is positive at all five points and
$f,g$ are nonzero.
\end{theorem}

This result concerns the displayed support. Theorem~\ref{thm:five-point}
allows every five-point integer support when $k$ is constant. Neither
theorem assumes a bound on the size or diameter of the support of $f$.

\subsection{Allocations from a common input}

Initially suppose that $f\ne0$ and that $g$ is positive on a five-point
set $F=\{p_1,\ldots,p_5\}$ and zero elsewhere. Write
$P=\|f\star g\|_1>0$. For each positive term of
\[
 N=\|f\star\widetilde{gk}\|_1
   =\sum_y\max_i k_i g_i f(y+p_i),
\]
select a maximizing index, resolving ties by index order. Let $E_i$
be the disjoint sets of outputs assigned to index $i$, and put
\[
 a_i=\frac{g_i}{P}\sum_{y\in E_i}f(y+p_i).
\]
Then $N/P=\sum_i k_i a_i$. Denominator outputs at $y+2p_i$ give
$0\le a_i\le1$. For distinct $i,j$, translate both $E_i$ and $E_j$
by $p_i+p_j$. They remain disjoint, and the corresponding denominator
witnesses give
\begin{equation}\label{weighted-pairs}
 \frac{g_j}{g_i}a_i+\frac{g_i}{g_j}a_j\le1,
 \qquad a_i a_j\le\frac14.
\end{equation}
Thus these constraints hold simultaneously for allocations obtained
from the same $f,g$, even though their selection depends on $k$.

\begin{lemma}\label{weighted-total}
Every allocation above satisfies
\[
 \sum_i a_i\le c:=\frac{1109}{450}.
\]
\end{lemma}
\begin{proof}
Put $R=\max_i g_i/\min_i g_i$ and $t=1/R$. We first show that
\begin{equation}\label{weighted-range}
 A:=\sum_i a_i\le F(t):=2+t-\frac{t^2}{2}.
\end{equation}
The pair at the largest and smallest values of $g$ satisfies
$ta+b/t\le1$. If $a,b\le v\le1$, filling the coordinate with
the smaller coefficient first gives
\[
 a+b\le\min\{2v,v+t-vt^2\}.
\]
When $M:=\max_i a_i\le1/2$, take $v=1/2$ and bound the other
three coordinates by $1/2$ to obtain \eqref{weighted-range}.

For $1/2\le M\le1$, all other coordinates are at most $1/(4M)$
by \eqref{weighted-pairs}. If the maximum is outside the extremal
pair, the preceding two-variable bound gives
\[
 A\le\min\{U(M),V(M)\},\qquad
 U(M)=M+\frac1M,\quad V(M)=M+\frac{3-t^2}{4M}+t.
\]
If the maximum belongs to the pair, assigning it to the smaller
coefficient can only increase the allowed partner. This gives the
same first bound and the second bound
$M+3/(4M)+t-Mt^2\le V(M)$.

The functions switch order at $M_0=(1+t^2)/(4t)\ge1/2$.
The function $U$ decreases on $[1/2,1]$, while $V$ is convex there.
Consequently it suffices to check $V(1/2)=F(t)$,
$V(1)\le F(t)$, and, if $M_0\le1$, $U(M_0)\le F(t)$.
The endpoint difference is $F(t)-V(1)=(1-t^2)/4$. At the switch,
\[
 4t(1+t^2)\bigl(F(t)-U(M_0)\bigr)
 =(1-t)^2(-2t^3-t^2+6t-1).
\]
Here $M_0\le1$ implies $t\ge2-\sqrt3>1/4$. On $[1/4,1]$,
the last factor is at least
$-3t^2+6t-1=2-3(1-t)^2\ge5/16$. This proves
\eqref{weighted-range}.

At each output, the selected unweighted product is at most
$\max_i g_i f(y+p_i)$. Therefore the energy estimate
\eqref{five-range-energy} also gives $A\le9R/5$. Combining the
two bounds yields
\[
 A\le\min\left\{\frac{9R}{5},\,
                  2+\frac1R-\frac1{2R^2}\right\}.
\]
For $R\le15/11$, the first expression is at most $27/11<c$.
For $R\ge15/11$, the second is decreasing and is at most $c$.
\end{proof}

For the kernel, write $S=\sum_i k_i$, $K=\max_i k_i$, and
$m=\min_i k_i$. Define
\[
 \delta=\frac52-c=\frac8{225},\qquad
 \alpha=\frac{c-\sqrt{c^2-4}}2,\qquad
 \beta=\frac1{4\alpha}.
\]
Then $1/2<\alpha<1$ and $\alpha+4\beta=c$.

\begin{lemma}\label{weighted-envelope}
The allocations satisfy
\[
 \sum_i k_i a_i\le\max\{U_c,V_c,C\},\qquad
 U_c=\frac S2-\delta m,\quad
 V_c=\alpha K+\beta(S-K),\quad C=\frac{S+3K}{4}.
\]
\end{lemma}
\begin{proof}
Use only $0\le a_i\le1$, $a_i a_j\le1/4$, and $\sum_i a_i\le c$.
These constraints are invariant under coordinate permutations, so
rearrangement permits a largest allocation $M$ at a largest kernel weight.
If $M\le1/2$, the deficit from five halves is at least $\delta$;
its weighted cost is at least $\delta m$, giving $U_c$.

If $1/2\le M\le\alpha$, the other four coordinates are at most
$1/(4M)$. Since $M+1/M\ge c$, subtracting the least possible cost
of the missing mass bounds the objective by
\[
 mc+(K-m)M+\frac{S-K-4m}{4M}.
\]
The coefficient $S-K-4m$ is nonnegative, so this expression is convex
in $M$. Its endpoint values are $U_c$ and $V_c$.
For $\alpha\le M\le1$, the convex bound
$KM+(S-K)/(4M)$ has endpoint values $V_c$ and $C$.
\end{proof}

\subsection{The corner comparison}

\begin{lemma}\label{weighted-corner}
If $F$ is a five-point Sidon set with $|4F|\ge32$, then
\[
 T_F(k)\ge\left(\frac{S+3K}{4}\right)^4.
\]
The inequality is strict when at least two kernel coordinates are positive.
\end{lemma}
\begin{proof}
The zero kernel is immediate. Normalize $K=1$, choose a position $p_i$
with weight one, and let $u_1,\ldots,u_4$ be the other weights.
Set $t=\sum_j u_j\in[0,4]$. Sidonicity makes the 15 outputs
$2p_i+(p_j+p_\ell)$, $j\le\ell$, distinct. Their total contribution
is at least
\[
 1+t+\frac{t^2+\sum_j u_j^2}{2}\ge1+t+\frac{5t^2}{8}.
\]
At least 17 other fourfold outputs each contribute at least $m^4$,
where $m\ge(t-3)_+$. Hence
\[
 T_F(k)-(1+t/4)^4
 \ge t^2 A(t)+17(t-3)_+^4,\qquad
 A(t)=\frac14-\frac t{16}-\frac{t^2}{256}.
\]
The function $A$ decreases on $[0,4]$. For $0<t\le13/4$ it is
at least $23/4096>0$. On the remaining intervals, the displayed
lower bound is bounded below respectively by
\[
\begin{array}{c|c}
 [13/4,10/3]&-25/1296+17/256=977/20736\\[2pt]
 [10/3,7/2]&-833/4096+17/81=2159/331776\\[2pt]
 [7/2,4]&-1+17/16=1/16.
\end{array}
\]
All are positive. When $t=0$, both sides equal one.
\end{proof}

For $F=\{0,2,7,8,11\}$, all 15 unordered pair sums are distinct,
and direct addition gives
\begin{equation}\label{weighted-fourfold-support}
 4F=\{0,2,4\}\cup\{6,7,\ldots,38\}\cup\{40,41,44\}.
\end{equation}
Thus $|4F|=39$ and Lemma~\ref{weighted-corner} applies.

\subsection{Rational certificates on the weight orders}

It remains to compare both $U_c$ and $V_c$ with $T_F(k)^{1/4}$.
Use the rational majorant
\[
 \overline V=aK+b(S-K),\qquad
 a=\frac{41}{80},\quad b=\frac{7027}{14400}=\frac{c-a}{4}.
\]
Indeed $a+1/a=8081/3280<c$, so $\alpha<a$. The function
$uK+(c-u)(S-K)/4$ has derivative $(5K-S)/4\ge0$;
therefore $V_c\le\overline V$.

Index the five positions of $F$ in increasing order. On a cone where
their weights have order $\pi$, write
\[
 k_{\pi(j)}=m+d_1+\cdots+d_{j-1}\quad(1\le j\le5),
 \qquad x=(m,d_1,d_2,d_3,d_4)\ge0.
\]
There are $5!=120$ orders, with ties allowed. The two linear targets
$U_c=q^U\cdot x$ and $\overline V=q^V\cdot x$ have vectors
\begin{align*}
 q^U&=(1109,900,675,450,225)/450,\\
 q^V&=(35488,28461,21434,14407,7380)/14400.
\end{align*}

Let $\mathcal M=\{\mu\in\mathbb N_0^5:|\mu|=4\}$ be the 70
unordered fourfold multisets. For $\mu\in\mathcal M$, put
\[
 z(\mu)=\sum_i p_i\mu_i,\qquad W_\mu(k)=\prod_i k_i^{\mu_i},
 \qquad \mathcal M_z=\{\mu:z(\mu)=z\}.
\]
There are 39 nonempty fibers $\mathcal M_z$ by
\eqref{weighted-fourfold-support}. These witness monomials have no
ordered-representation multiplicity: the operation is max-convolution.

For each order $\pi$ and each target $q$, the certificate supplies 70
nonnegative integers $n_\mu$, with common denominator $D=10^9$,
satisfying
\begin{equation}\label{weighted-certificate}
 \sum_{\mu\in\mathcal M_z}n_\mu=D\quad(z\in4F),\qquad
 \sum_\mu\frac{n_\mu}{D}W_\mu(k(x))-(q\cdot x)^4
 \text{ has nonnegative coefficients.}
\end{equation}
The supplied verifier checks these conditions for all 240 pairs
$(\pi,q)$. Its arithmetic uses integers throughout each coefficient
comparison. If $q=v/Q$, the comparison for a degree-four exponent
$\gamma$ is
\[
 Q^4\sum_\mu n_\mu[x^\gamma]W_\mu(k(x))
       -D[x^\gamma](v\cdot x)^4\ge0.
\]
The program reconstructs every witness coefficient by enumerating the
variable choices in its four slots, and expands the target by repeated
polynomial multiplication. It checks the fiber masses, all 70 coefficients,
and both targets for every order. The coefficient inequalities are checked
independently of the solver's floating-point values.

For every $x\ge0$, \eqref{weighted-certificate} implies
\[
 T_F(k)=\sum_z\max_{\mu\in\mathcal M_z}W_\mu(k)
 \ge\sum_\mu\frac{n_\mu}{D}W_\mu(k)
 \ge(q\cdot x)^4.
\]
Every witness has $m^4$ coefficient one. Fiber normalization therefore
makes that coefficient 39 in the middle polynomial, whereas both target
fourth powers have coefficient $c^4<3721/100<39$. Thus both comparisons
are strict when $m>0$.

\begin{proof}[Proof of Theorem~\ref{thm:five-point-weighted}]
For positive $k,g$ on $F$ and nonzero $f$, combine the allocation
envelope with the two certificate comparisons and the corner lemma.
Each of $U_c,V_c,C$ is strictly below $T_F(k)^{1/4}$, so the
reflected inequality is strict.

For fixed $f$, all relevant norms are finite sums of maxima of finitely
many coordinate products, and hence continuous in $k,g$. Positive
approximation gives the non-strict inequality when some coordinates
vanish. When $k>0$, the upper envelope depends only on $k$ and stays
strictly below $T_F(k)^{1/4}$; taking the same limit in the envelope
therefore preserves strictness for nonzero $g$ supported on $F$.

For general $g$, restrict it to $g_F=g1_F$. The numerator is unchanged,
while $\|f\star g_F\|_1\le\|f\star g\|_1$. If $g_F\ne0$, apply
the preceding result. If $g_F=0$, the numerator is zero; when $k>0$
and $f,g\ne0$, the right side is positive. The zero-function cases
of the non-strict statement are immediate.
\end{proof}

The certificate files and standard-library checker accompany the paper.
From this paper's source directory, run
\begin{verbatim}
python3 -B verify_five_point.py
\end{verbatim}
The command checks all 240 certificates and corruption controls. It
requires no numerical optimizer or external Python package. The finite
certificates establish the coefficient comparisons; the preceding hand
argument supplies their connection to the reflected inequality.


\paragraph{Review and assistance.}
This manuscript was developed with AI assistance. Separate AI agents checked
the mathematical arguments and their transcription. These checks do not
constitute human peer review or journal acceptance.

\begin{thebibliography}{9}
\small\raggedright
\setlength{\itemsep}{3pt}
\bibitem{ghr}
K.~Gyarmati, F.~Hennecart and I.~Z.~Ruzsa,
\emph{Sums and differences of finite sets},
Functiones et Approximatio Commentarii Mathematici \textbf{37}(1) (2007),
175--186,
\href{https://doi.org/10.7169/facm/1229618749}{doi:10.7169/facm/1229618749}.
\bibitem{mrsz}
D.~Matolcsi, I.~Z.~Ruzsa, G.~Shakan and D.~Zhelezov,
\emph{An analytic approach to cardinalities of sumsets},
\href{https://arxiv.org/abs/2003.04075v1}{arXiv:2003.04075v1} (2020).
\bibitem{bikm}
L.~Becker, P.~Ivanisvili, D.~Krachun and J.~Madrid,
\emph{Discrete Brunn--Minkowski inequality for subsets of the cube},
\href{https://arxiv.org/abs/2404.04486v2}{arXiv:2404.04486v2} (2024).
\bibitem{hi}
J.~Hosle and P.~Ivanisvili,
\emph{Geometric Block Exponents and a Uniform Mixed-Alphabet Sumset Inequality},
\href{https://arxiv.org/abs/2606.25350v1}{arXiv:2606.25350v1} (2026).
\bibitem{cs}
A.~Cloninger and S.~Steinerberger,
\emph{On suprema of autoconvolutions with an application to Sidon sets},
\href{https://arxiv.org/abs/1403.7988v3}{arXiv:1403.7988v3} (2016).
Published in Proceedings of the American Mathematical Society
\textbf{145}(8) (2017), 3191--3200.
\bibitem{ix}
P.~Ivanisvili and X.~Xie,
\emph{Grokability in five inequalities},
\href{https://arxiv.org/abs/2605.05193v1}{arXiv:2605.05193v1} (2026).
\bibitem{blt}
B.~Bollob\'as, I.~Leader and M.~Tiba,
\emph{A strengthening of Freiman's $3k-4$ theorem},
\href{https://arxiv.org/abs/2204.09816v1}{arXiv:2204.09816v1} (2022).
\bibitem{gmrsz}
B.~Green, D.~Matolcsi, I.~Z.~Ruzsa, G.~Shakan and D.~Zhelezov,
\emph{A Weighted Pr\'ekopa--Leindler inequality and sumsets with quasicubes},
\href{https://arxiv.org/abs/2003.04077v1}{arXiv:2003.04077v1} (2020).
Published in A.~Avila, M.~T.~Rassias and Y.~Sinai (eds.),
\emph{Analysis at Large}, Springer, 2022, pp.~125--129.
\end{thebibliography}
\end{document}
